\chapter{K\"ahler Manifolds and the Hodge Decomposition}\label{kahler}
\chapterstatus{complete}

\section{Introduction}\label{kahler:section-introduction}

A K\"ahler metric is a Hermitian metric whose fundamental $2$-form is closed. This single condition makes the Riemannian and the complex
structure compatible to second order (Proposition~\ref{kahler:proposition-kahler-equivalent}), it holds on every projective manifold
(Corollary~\ref{kahler:corollary-projective-kahler}), and it is exactly what the Hodge decomposition of Section~\ref{kahler:section-hodge-decomposition}
needs. This chapter develops K\"ahler metrics and the pointwise linear algebra of the Lefschetz operator
(Theorem~\ref{kahler:theorem-sl2-forms}), which becomes the hard Lefschetz theorem in Section~\ref{kahler:section-lefschetz-theorems}. References:
\cite[Chapter 0]{griffiths-harris}, \cite[Chapters 1 and 3]{huybrechts-complex}, \cite[Chapters 3 and 6]{voisin-hodge-1}.

\noindent\textbf{Prerequisites.} Section~\ref{complex-manifolds:section-dolbeault} (Dolbeault Cohomology), Chapter~\ref{riemannian-geometry} (Riemannian
Geometry) and Chapter~\ref{representation-theory} (Representations of Groups and Lie Algebras).

Throughout, $X$ is a complex manifold of complex dimension $n$ and $J$ is its almost complex structure (multiplication by $i$ on tangent
vectors). Forms are complex valued unless said otherwise, and $\Lambda^{p,q}$ denotes forms of type $(p,q)$.

\section{Hermitian metrics and the fundamental form}\label{kahler:section-hermitian}

\begin{definition}\label{kahler:definition-hermitian}
A \emph{Hermitian metric} on $X$ is a Riemannian metric $g$ with $g(JU, JV) = g(U, V)$ for all tangent vectors $U, V$. Its
\emph{fundamental form} is the real $2$-form $\omega(U, V) = g(JU, V)$.
\end{definition}

\begin{proposition}\label{kahler:proposition-fundamental-form-coordinates}
Let $g$ be a Hermitian metric and extend it complex bilinearly to complexified tangent vectors. In holomorphic coordinates
$z^j = x^j + iy^j$ put $h_{j\bar k} = 2g(\partial/\partial z^j, \partial/\partial\bar z^k)$. Then $(h_{j\bar k})$ is a positive definite
Hermitian matrix at every point and
\[
\omega = \frac{i}{2}\sum_{j,k}h_{j\bar k}\,dz^j \wedge d\bar z^k,
\]
a real form of type $(1,1)$. Conversely, every real $(1,1)$-form $\frac{i}{2}\sum h_{j\bar k}\,dz^j \wedge d\bar z^k$ whose matrix
$(h_{j\bar k})$ is positive definite (such a form is called \emph{positive}) is the fundamental form of exactly one Hermitian metric. For the
standard metric of $\CC^n$ one gets $\omega_0 = \frac{i}{2}\sum_jdz^j \wedge d\bar z^j = \sum_jdx^j \wedge dy^j$.
\end{proposition}

\begin{proof}
The $J$-invariance of $g$ says that $g$ pairs the $+i$- and $-i$-eigenspaces of $J$ with each other and each of them with itself
trivially: $g(\partial_j, \partial_k) = g(J\partial_j, J\partial_k) = i^2g(\partial_j, \partial_k)$ forces $g(\partial_j, \partial_k) = 0$,
where $\partial_j = \partial/\partial z^j$, and similarly for $\bar\partial_j = \partial/\partial\bar z^j$. So
$g = \frac12\sum h_{j\bar k}(dz^j \otimes d\bar z^k + d\bar z^k \otimes dz^j)$. Since $g$ is real, $\overline{h_{j\bar k}} = h_{k\bar j}$,
and for $U = \sum u^j\partial_j + \overline{u^j}\bar\partial_j$ real, $g(U, U) = \sum h_{j\bar k}u^j\overline{u^k}$, so $(h_{j\bar k})$ is
positive definite. Since $J\partial_j = i\partial_j$ and $J\bar\partial_k = -i\bar\partial_k$,
$\omega(\partial_j, \bar\partial_k) = g(i\partial_j, \bar\partial_k) = \frac{i}{2}h_{j\bar k}$,
$\omega(\bar\partial_k, \partial_j) = -\frac{i}{2}h_{j\bar k}$ and $\omega$ vanishes on pairs of the same type; this is the displayed
formula. Conversely, a positive $(1,1)$-form determines $h_{j\bar k}$ and hence $g$ by the same formulas, and $g$ is then a $J$-invariant
Riemannian metric. For $\CC^n$, $g(\partial_j, \bar\partial_k) = \frac14(g(\partial_{x^j}, \partial_{x^k}) + g(\partial_{y^j}, \partial_{y^k})) = \frac12\delta_{jk}$, and
$\frac{i}{2}dz \wedge d\bar z = dx \wedge dy$.
\end{proof}

\begin{remark}\label{kahler:remark-two-of-three}
At every point $x$, the metric $g_x$, the fundamental form $\omega_x$ and the complex structure $J_x$ form a compatible triple on $T_xX$ in the sense of
Proposition~\ref{linear-algebra:proposition-two-of-three}, so any two of them determine the third.
\begin{enumerate}
\item $g$ and $J$ give $\omega(U, V) = g(JU, V)$: this is the definition.
\item $\omega$ and $J$ give $g(U, V) = \omega(U, JV)$, provided $\omega(JU, JV) = \omega(U, V)$ and $\omega(U, JU) > 0$ for $U \neq 0$. These two conditions say, in coordinates, that $\omega$ is a real
$(1,1)$-form with positive definite matrix $(h_{j\bar k})$, which is the converse part of Proposition~\ref{kahler:proposition-fundamental-form-coordinates}. So a Hermitian manifold can equally be
described by $(X, \omega)$ with $\omega$ a positive $(1,1)$-form, which is how K\"ahler metrics are usually given.
\item $g$ and $\omega$ give $J$ through $\omega(U, V) = g(JU, V)$. Here the pair is restricted twice: the endomorphism $A$ defined by $\omega(U, V) = g(AU, V)$ must satisfy $A^2 = -1$ for $A$ to be an almost
complex structure, and the almost complex structure must moreover be integrable (Theorem~\ref{complex-manifolds:theorem-newlander-nirenberg}) to come from a complex manifold.
\end{enumerate}
Correspondingly, a linear map of $T_xX$ preserving two of the three structures preserves the third, and the maps preserving all three form the unitary group $\mathrm{U}(n)$. The analogous ``two out of
three'' statement one level up concerns the K\"ahler condition: for a Hermitian metric, $d\omega = 0$ holds if and only if $J$ is parallel for the Levi-Civita connection of $g$
(Remark~\ref{kahler:remark-levi-civita}), and then $g$, $\omega$ and $J$ are all parallel. Dropping the integrability of $J$ and keeping $d\omega = 0$ gives the \emph{almost K\"ahler} manifolds, which are the symplectic manifolds with a
compatible almost complex structure (Chapter~\ref{symplectic-geometry}).
\end{remark}

\begin{proposition}\label{kahler:proposition-volume}
For a Hermitian metric, $\omega^n/n! = \mathrm{vol}_g$ is the Riemannian volume form for the canonical orientation. The restriction of $\omega$
to a complex submanifold $Y$ is the fundamental form of the restricted metric. Hence for a nonempty compact complex submanifold, or analytic
subvariety, $Z$ of pure dimension $k$,
\[
\int_Z\frac{\omega^k}{k!} = \mathrm{vol}(Z_{\mathrm{reg}}) > 0 .
\]
\end{proposition}

\begin{proof}
At a point choose, by Gram--Schmidt for the Hermitian form $(h_{j\bar k})$, a linear change of coordinates making $h_{j\bar k} = \delta_{jk}$
there. Then $\omega = \sum_jdx^j \wedge dy^j$ at that point, so
$\omega^n = n!\,dx^1 \wedge dy^1 \wedge \cdots \wedge dx^n \wedge dy^n$, which is the volume form of the orthonormal basis
$\partial_{x^1}, \partial_{y^1}, \ldots$ in the canonical orientation. The tangent spaces of $Y$ are $J$-stable, so the restricted metric is
Hermitian, and $\omega(U, V) = g(JU, V)$ for $U, V$ tangent to $Y$ is its fundamental form. The formula for $Z$ follows by applying the
first statement to $Z_{\mathrm{reg}}$, which is a nonempty open complex manifold of dimension $k$; the integral converges by
Proposition~\ref{complex-manifolds:proposition-subvariety-class}.
\end{proof}

\section{The K\"ahler condition}\label{kahler:section-kahler}

\begin{definition}\label{kahler:definition-kahler}
A Hermitian metric is \emph{K\"ahler} if its fundamental form is closed, $d\omega = 0$; then $\omega$ is called a \emph{K\"ahler form}. A
complex manifold is \emph{K\"ahler} if it admits a K\"ahler metric, and the class $[\omega] \in H^2(X; \RR)$ is then a
\emph{K\"ahler class}.
\end{definition}

\begin{proposition}\label{kahler:proposition-kahler-equivalent}
For a Hermitian metric on $X$ the following are equivalent:
\begin{enumerate}
\item $d\omega = 0$;
\item around every point $p$ there are holomorphic coordinates centred at $p$ with $h_{j\bar k} = \delta_{jk} + O(|z|^2)$ (the metric
\emph{osculates to order two} to the flat metric);
\item locally $\omega = i\partial\bar\partial\varphi$ for a smooth real function $\varphi$, a \emph{local K\"ahler potential}.
\end{enumerate}
\end{proposition}

\begin{proof}
(2)$\Rightarrow$(1): in the coordinates of (2), the coefficients of $d\omega$ at $p$ are first derivatives of the $h_{j\bar k}$ at $0$, which
vanish; since $p$ is arbitrary, $d\omega = 0$. (3)$\Rightarrow$(1): $d(i\partial\bar\partial\varphi) = i(\partial + \bar\partial)\partial\bar\partial\varphi = 0$ by
Proposition~\ref{complex-manifolds:proposition-dbar}.

(1)$\Rightarrow$(2): after a linear change of coordinates, $h_{j\bar k}(0) = \delta_{jk}$. By Taylor's theorem
\[
h_{j\bar k} = \delta_{jk} + \sum_l\left(a_{jkl}z^l + b_{jkl}\bar z^l\right) + O(|z|^2),
\]
and $\overline{h_{j\bar k}} = h_{k\bar j}$ gives $b_{jkl} = \overline{a_{kjl}}$. The $(2,1)$-component of $d\omega$ is
$\partial\omega = \frac{i}{2}\sum\partial_lh_{j\bar k}\,dz^l \wedge dz^j \wedge d\bar z^k$, and its vanishing at $0$ says that
$a_{jkl} = a_{lkj}$. Define $w^k = z^k + \frac12\sum_{j,l}a_{jkl}z^jz^l$. Since $dw^k = dz^k$ at $0$, the map $z \mapsto w$ is a diffeomorphism near $0$
(Theorem~\ref{real-analysis:theorem-inverse-function}), and its inverse is holomorphic because its derivative, the inverse of a complex
linear map, is complex linear; so the $w^k$ are holomorphic coordinates near $0$. By the symmetry $a_{jkl} = a_{lkj}$,
$dw^k = dz^k + \sum_{j,l}a_{jkl}z^l\,dz^j$. Hence
\[
\sum_kdw^k \wedge d\bar w^k = \sum_{j,k}\Big(\delta_{jk} + \sum_la_{jkl}z^l + \sum_l\overline{a_{kjl}}\,\bar z^l\Big)dz^j \wedge d\bar z^k + O(|z|^2),
\]
so $\omega - \frac{i}{2}\sum_kdw^k \wedge d\bar w^k$ has coefficients $O(|z|^2)$ in the basis $dz^j \wedge d\bar z^k$, hence also in the
basis $dw^j \wedge d\bar w^k$, which differs from it by a smooth change of basis, and $|z|$ and $|w|$ are comparable near $0$. This is (2).

(1)$\Rightarrow$(3): on a polydisc $\Delta$ around a point, $\omega = d\eta$ for a real $1$-form $\eta$ by the Poincar\'e lemma
(Theorem~\ref{de-rham:theorem-homotopy-invariance}). Write $\eta = \eta^{1,0} + \eta^{0,1}$ with $\eta^{1,0} = \overline{\eta^{0,1}}$.
The $(0,2)$-component of $d\eta = \omega$ is $\bar\partial\eta^{0,1} = 0$, so on a smaller polydisc $\eta^{0,1} = \bar\partial f$ for a
smooth complex function $f$ by Lemma~\ref{complex-manifolds:lemma-dbar-poincare-pq}, and then $\eta^{1,0} = \partial\bar f$. The $(1,1)$-component
gives
\[
\omega = \partial\eta^{0,1} + \bar\partial\eta^{1,0} = \partial\bar\partial f + \bar\partial\partial\bar f = \partial\bar\partial(f - \bar f),
\]
so $\varphi = -i(f - \bar f) = 2\operatorname{Im}f$ is real and $i\partial\bar\partial\varphi = \omega$.
\end{proof}

\begin{remark}\label{kahler:remark-levi-civita}
The K\"ahler condition has several other equivalent forms: $J$ is parallel for the Levi-Civita connection; the Levi-Civita connection of $g$,
transported to $T^{1,0}X$, is the Chern connection of the Hermitian holomorphic bundle $T^{1,0}X$
(Theorem~\ref{complex-manifolds:theorem-chern-connection}); the holonomy group of $g$ is contained in $U(n)$ \cite[Chapter 4]{huybrechts-complex}.
Condition (2) is the useful one in practice: any first-order identity between operators built from $g$ and $J$ that holds on $\CC^n$ with
the flat metric holds on every K\"ahler manifold, since it can be checked at the centre of such coordinates. The K\"ahler identities of
Section~\ref{kahler:section-hodge-decomposition} are proved this way.
\end{remark}

\section{Examples}\label{kahler:section-examples}

\begin{proposition}[Fubini--Study metric]\label{kahler:proposition-fubini-study}
Let $\sigma$ be a holomorphic section of the projection $\CC^{n+1} \setminus \{0\} \to \CC\PP^n$ over an open set $U$. The form
\[
\omega_{FS} = \frac{i}{2\pi}\partial\bar\partial\log|\sigma|^2
\]
does not depend on $\sigma$, and so defines a real $(1,1)$-form on $\CC\PP^n$. It is a K\"ahler form, invariant under the action of
$U(n+1)$, and $[\omega_{FS}] = c_1(\mathcal{O}(1))$, the positive generator of $H^2(\CC\PP^n; \ZZ)$.
\end{proposition}

\begin{proof}
Another section is $\lambda\sigma$ with $\lambda$ holomorphic and nowhere zero on $U$ (shrinking $U$ so that $\log\lambda$ exists), and
$\log|\lambda\sigma|^2 = \log|\sigma|^2 + \log\lambda + \overline{\log\lambda}$. Since $\bar\partial\log\lambda = 0$ and
$\partial\overline{\log\lambda} = 0$, $\partial\bar\partial$ kills the last two terms, so $\omega_{FS}$ is well defined; it is real because
$\overline{i\partial\bar\partial u} = -i\bar\partial\partial u = i\partial\bar\partial u$ for real $u$, and closed by
Proposition~\ref{kahler:proposition-kahler-equivalent}. For $A \in U(n+1)$, $A \circ \sigma$ is a section over $A(U)$ with
$|A\sigma| = |\sigma|$, so $\omega_{FS}$ is $U(n+1)$-invariant. As $U(n+1)$ acts transitively on $\CC\PP^n$, positivity need only be checked
at $[1 : 0 : \cdots : 0]$. With the section $\sigma(w) = (1, w_1, \ldots, w_n)$ in affine coordinates,
$\log|\sigma|^2 = \log(1 + |w|^2) = |w|^2 + O(|w|^4)$, so at $w = 0$
\[
\omega_{FS} = \frac{i}{2\pi}\sum_jdw^j \wedge d\bar w^j = \frac1\pi\sum_jdu^j \wedge dv^j \qquad (w^j = u^j + iv^j),
\]
which is positive. Finally, $\sigma$ is a nowhere-vanishing section of $\mathcal{O}(-1) \subseteq \CC\PP^n \times \CC^{n+1}$, and the dual
frame $\sigma^*$ of $\mathcal{O}(1)$ has norm $|\sigma^*|^2 = |\sigma|^{-2} = e^{-\varphi}$ for the dual of the standard metric, with
$\varphi = \log|\sigma|^2$. By Theorem~\ref{complex-manifolds:theorem-chern-connection}, $c_1(\mathcal{O}(1)) = [\frac{i}{2\pi}\partial\bar\partial\varphi] = [\omega_{FS}]$.
This is the positive generator of $H^2(\CC\PP^n; \ZZ)$ by Example~\ref{complex-manifolds:example-fubini-study}.
\end{proof}

\begin{example}\label{kahler:example-fubini-study-explicit}
\emph{The Fubini--Study form in coordinates.} On $U_0 = \{[1 : w_1 : \cdots : w_n]\} \cong \CC^n$, with $\sigma(w) = (1, w)$ and $|w|^2 = \sum_j|w_j|^2$,
$\bar\partial\log(1 + |w|^2) = \sum_kw_k\,d\bar w_k/(1 + |w|^2)$, and applying $\partial$,
\[
\omega_{FS} = \frac{i}{2\pi}\Bigl(\frac{\sum_jdw_j \wedge d\bar w_j}{1 + |w|^2} - \frac{(\sum_j\bar w_j\,dw_j) \wedge (\sum_kw_k\,d\bar w_k)}{(1 + |w|^2)^2}\Bigr)
= \frac{i}{2}\sum_{j,k}\frac{h_{j\bar k}}{\pi}\,dw_j \wedge d\bar w_k, \quad h_{j\bar k} = \frac{(1 + |w|^2)\delta_{jk} - \bar w_jw_k}{(1 + |w|^2)^2}.
\]
This shows positivity at every point directly, without the homogeneity argument: for $\xi \in \CC^n$, by the Cauchy--Schwarz inequality $|\sum_kw_k\xi_k|^2 \leq |w|^2|\xi|^2$,
\[
\sum_{j,k}\bar\xi_jh_{j\bar k}\xi_k = \frac{(1 + |w|^2)|\xi|^2 - |\sum_kw_k\xi_k|^2}{(1 + |w|^2)^2} \geq \frac{|\xi|^2}{(1 + |w|^2)^2} > 0 .
\]
For $n = 1$, with $w = u + iv$ and $dw \wedge d\bar w = -2i\,du \wedge dv$, $\omega_{FS} = \frac1\pi\frac{du \wedge dv}{(1 + |w|^2)^2}$, and the area of $\CC\PP^1$ is
$\int_{U_0}\omega_{FS} = \frac1\pi\int_0^{2\pi}\int_0^\infty\frac{r\,dr\,d\theta}{(1 + r^2)^2} = 1$, the complement of $U_0$ being a point. This agrees with $[\omega_{FS}] = c_1(\mathcal{O}(1))$, whose degree is
$1$; in general $\int_{\CC\PP^n}\omega_{FS}^n = 1$, so the Riemannian volume of $\CC\PP^n$ for this normalisation is $\frac1{n!}$, as $\omega^n/n!$ is the volume form. On $\CC\PP^1$ the metric of
$\pi\omega_{FS}$ is invariant under the transitive isometric action of $U(2)$, so its Gaussian curvature $K$ is constant; its area is $\pi$, so $K\pi = 4\pi$ by the Gauss--Bonnet theorem
(Theorem~\ref{riemannian-geometry:theorem-global-gauss-bonnet}), $K = 4$, and $\CC\PP^1$ with this metric is isometric to the round sphere of radius $\frac12$
(Theorem~\ref{riemannian-geometry:theorem-killing-hopf}).
\end{example}

\begin{corollary}\label{kahler:corollary-projective-kahler}
A complex submanifold of a K\"ahler manifold is K\"ahler with the restricted metric. In particular every \emph{projective manifold}, that is,
every closed complex submanifold $X \subseteq \CC\PP^N$, is K\"ahler, and its K\"ahler class $[\omega_{FS}|_X]$ is the restriction of the
hyperplane class $c_1(\mathcal{O}(1))$.
\end{corollary}

\begin{proof}
By Proposition~\ref{kahler:proposition-volume}, the fundamental form of the restricted metric is the restriction of $\omega$, and restriction
of forms commutes with $d$.
\end{proof}

\begin{example}\label{kahler:example-kahler}
\begin{enumerate}
\item $\CC^n$ with $\omega_0$, and complex tori $\CC^n/\Lambda$ with the form induced by $\omega_0$, which is invariant under translations. Products
of K\"ahler manifolds are K\"ahler, with $\omega = \mathrm{pr}_1^*\omega_1 + \mathrm{pr}_2^*\omega_2$.
\item Every Hermitian metric on a Riemann surface is K\"ahler, since $d\omega$ is a $3$-form on a real surface.
\item Every smooth projective variety is K\"ahler (Corollary~\ref{kahler:corollary-projective-kahler}). The blow-up of a compact K\"ahler
manifold along a compact complex submanifold is again K\"ahler \cite[Chapter 3]{voisin-hodge-1}.
\item Not every compact K\"ahler manifold is projective: a complex torus $\CC^n/\Lambda$ with $n \geq 2$ and $\Lambda$ general has no
nonconstant meromorphic functions, hence no embedding into projective space \cite{birkenhake-lange}; see also Chapter~\ref{algebraic-groups}. Voisin even showed that there
are compact K\"ahler manifolds that are not homotopy equivalent to any projective manifold \cite{voisin-kodaira}.
\end{enumerate}
\end{example}

\begin{proposition}\label{kahler:proposition-kahler-class-powers}
Let $X$ be a compact K\"ahler manifold with K\"ahler form $\omega$. Then $[\omega]^k \neq 0$ in $H^{2k}(X; \RR)$ for $0 \leq k \leq n$; in
particular $b_{2k}(X) \geq 1$. The cycle class (Proposition~\ref{complex-manifolds:proposition-subvariety-class}) of a nonempty analytic
subvariety $Z$ of pure dimension $k$ is nonzero in $H^{2n-2k}(X; \RR)$.
\end{proposition}

\begin{proof}
By Proposition~\ref{kahler:proposition-volume}, $\int_X\omega^k \wedge \omega^{n-k} = n!\,\mathrm{vol}(X) > 0$. If $\omega^k = d\eta$ were
exact, then $\omega^n = d(\eta \wedge \omega^{n-k})$ would be exact and its integral would vanish by Stokes' theorem. For $Z$, the class
$[Z]$ is represented by the current of integration over $Z_{\mathrm{reg}}$, whose value on the closed form $\omega^k$ is
$\int_Z\omega^k = k!\,\mathrm{vol}(Z_{\mathrm{reg}}) > 0$; since a closed current that is exact vanishes on closed forms, $[Z] \neq 0$
(Theorem~\ref{currents:theorem-currents}).
\end{proof}

\begin{lemma}[Wirtinger inequality for $2$-planes]\label{kahler:lemma-wirtinger}
Let $g$ be a Hermitian metric with fundamental form $\omega$, and let $u$, $v$ be orthonormal tangent vectors at a point. Then $|\omega(u, v)| \leq 1$, with equality if and only if $v = \pm Ju$, that is, if and only if
the real $2$-plane spanned by $u$ and $v$ is a complex line.
\end{lemma}

\begin{proof}
$\omega(u, v) = g(Ju, v)$, and $|Ju| = |u| = 1$ because $g$ is $J$-invariant. By the Cauchy--Schwarz inequality $|g(Ju, v)| \leq |Ju||v| = 1$, with equality if and only if $v = \pm Ju$. A real plane containing $u$
and $Ju$ is $J$-stable, hence a complex line.
\end{proof}

\begin{proposition}\label{kahler:proposition-area-minimising}
Let $(X, \omega)$ be a K\"ahler manifold and $C \subseteq X$ a compact complex curve (a compact complex submanifold of dimension $1$). If $\Sigma$ is a compact oriented surface and $f \colon \Sigma \to X$ an immersion
with $f_*[\Sigma] = [C]$ in $H_2(X; \RR)$, then $\mathrm{area}(f) \geq \mathrm{area}(C) = \int_C\omega$, with equality if $f(\Sigma)$ is a complex curve with its complex orientation. So holomorphic curves minimise area in
their homology class.
\end{proposition}

\begin{proof}
At a point of $\Sigma$ choose an oriented orthonormal basis $u$, $v$ of $T\Sigma$ for the pulled-back metric; the area form of $f^*g$ takes the value $1$ on $(u, v)$, while $f^*\omega$ takes the value
$\omega(df\,u, df\,v) \leq 1$ by Lemma~\ref{kahler:lemma-wirtinger}. So $f^*\omega \leq dA$ pointwise, and
\[
\mathrm{area}(f) = \int_\Sigma dA \geq \int_\Sigma f^*\omega = \int_C\omega = \mathrm{area}(C),
\]
where the middle equality holds because $\omega$ is closed and $f_*[\Sigma] = [C]$ (Theorem~\ref{de-rham:theorem-de-rham}), and the last is Proposition~\ref{kahler:proposition-volume} with $k = 1$. If $f(\Sigma)$ is a
complex curve with the complex orientation, equality holds pointwise by the equality case of the lemma.
\end{proof}

\begin{remark}\label{kahler:remark-calibrations}
The same argument, with the Wirtinger inequality $\frac{\omega^k}{k!} \leq \mathrm{vol}$ on real $2k$-planes, shows that compact complex submanifolds of any dimension minimise volume in their homology class
\cite[Chapter 0, Section 2]{griffiths-harris}. A closed form that is bounded by the volume on every plane, with equality on a class of submanifolds, is called a \emph{calibration}; the K\"ahler form is the first
example, and the holomorphic volume form of a Calabi--Yau manifold calibrates the \emph{special Lagrangian} submanifolds \cite{joyce-special-holonomy}.
\end{remark}

\begin{counterexample}\label{kahler:counterexample-hopf}
The Hopf surface $X = (\CC^2 \setminus \{0\})/\langle z \mapsto 2z\rangle$ is a compact complex manifold diffeomorphic to $S^1 \times S^3$
(Exercise~\ref{complex-manifolds:exercise-hopf-surface}), and it is not K\"ahler: $H^2(X; \RR) = 0$, whereas a compact K\"ahler manifold of positive
dimension has $H^2 \neq 0$ by Proposition~\ref{kahler:proposition-kahler-class-powers}. Another obstruction, that $b_1(X) = 1$ is odd, is
explained in Section~\ref{kahler:section-hodge-decomposition}.
\end{counterexample}

\section{Linear algebra of the Lefschetz operator}\label{kahler:section-lefschetz-linear-algebra}

\begin{definition}\label{kahler:definition-lefschetz-operators}
Let $V$ be a real vector space of dimension $2n$ with a complex structure $J$ and a $J$-invariant inner product $g$, with fundamental form
$\omega$. Let $\Lambda^\bullet = \Lambda^\bullet V^* \otimes_\RR \CC$ with the Hermitian inner product induced by $g$, for which the wedge
products of an orthonormal basis of $V^*$ form an orthonormal basis. Define:
\begin{itemize}
\item the \emph{Lefschetz operator} $L(\alpha) = \omega \wedge \alpha$, of degree $+2$;
\item its adjoint $\Lambda = L^*$, of degree $-2$;
\item the \emph{counting operator} $H = \sum_k(k - n)\pi_k$, where $\pi_k$ is the projection onto $\Lambda^k$.
\end{itemize}
On a Hermitian manifold these operators act pointwise on forms. A form $\alpha$ with $\Lambda\alpha = 0$ is called \emph{primitive}, and
$P^k$ denotes the primitive forms of degree $k$.
\end{definition}

\begin{theorem}\label{kahler:theorem-sl2-forms}
The operators satisfy
\[
[L, \Lambda] = H, \qquad [H, L] = 2L, \qquad [H, \Lambda] = -2\Lambda,
\]
so $e = L$, $f = \Lambda$, $h = H$ define a representation of $\mathfrak{sl}_2$ on $\Lambda^\bullet$ in which $\Lambda^k$ is the weight space of
weight $k - n$. Consequently:
\begin{enumerate}
\item for $0 \leq k \leq n$, $L^{n-k} \colon \Lambda^k \to \Lambda^{2n-k}$ is an isomorphism;
\item $P^k = 0$ for $k > n$, and every $\alpha \in \Lambda^k$ can be written uniquely as $\alpha = \sum_{j \geq 0}L^j\alpha_j$ with
$\alpha_j \in P^{k-2j}$ (the \emph{Lefschetz} or \emph{primitive decomposition});
\item for $k \leq n$, a form $\alpha \in \Lambda^k$ is primitive if and only if $L^{n-k+1}\alpha = 0$;
\item $L$, $\Lambda$ map $\Lambda^{p,q}$ to $\Lambda^{p+1,q+1}$ and $\Lambda^{p-1,q-1}$, and if $\alpha$ has type $(p,q)$ then each
$\alpha_j$ has type $(p - j, q - j)$.
\end{enumerate}
\end{theorem}

\begin{proof}
The last two commutators hold because $L$ raises the degree by $2$ and $\Lambda$ lowers it by $2$: for $\alpha \in \Lambda^k$,
$HL\alpha - LH\alpha = (k + 2 - n)L\alpha - (k - n)L\alpha = 2L\alpha$, and similarly for $\Lambda$.

For $[L, \Lambda] = H$, choose an orthonormal basis $e_1, Je_1, \ldots, e_n, Je_n$ of $V$ (possible by induction: $e_1^\perp$ is
$J$-stable when $e_1$ is a unit vector and $e_1 \perp Je_1$, as $g(e_1, Je_1) = g(Je_1, J^2e_1) = -g(Je_1, e_1)$). Let
$x^1, y^1, \ldots, x^n, y^n$ be the dual basis, so that $\omega = \sum_j\omega_j$ with $\omega_j = x^j \wedge y^j$. The algebra
$\Lambda^\bullet$ is the graded tensor product of the four-dimensional algebras $A_j = \Lambda^\bullet\langle x^j, y^j\rangle$, with the
product inner product; $L = \sum_jL_j$ where $L_j$ is multiplication by $\omega_j$ in the $j$-th factor, and so $\Lambda = \sum_j\Lambda_j$
with $\Lambda_j$ the adjoint of $L_j$ acting in the $j$-th factor. Since $L_j$ and $\Lambda_j$ have even degree, operators in different
factors commute, and
\[
[L, \Lambda] = \sum_{j,k}[L_j, \Lambda_k] = \sum_j[L_j, \Lambda_j].
\]
In $A_j$, with orthonormal basis $1, x^j, y^j, x^j \wedge y^j$, the operator $L_j$ maps $1 \mapsto x^j \wedge y^j$ and kills the other basis
vectors, and $\Lambda_j$ maps $x^j \wedge y^j \mapsto 1$ and kills the others. So $[L_j, \Lambda_j]$ is $-1$ on $1$, $0$ on $x^j, y^j$ and
$+1$ on $x^j \wedge y^j$: it is multiplication by $(\deg_j - 1)$, where $\deg_j$ is the degree in the $j$-th factor. On a product of basis
elements of total degree $k = \sum_j\deg_j$, $[L, \Lambda]$ is therefore multiplication by $\sum_j(\deg_j - 1) = k - n$, which is $H$.

The consequences are statements about finite-dimensional $\mathfrak{sl}_2$-representations. The map $(e, f, h) \mapsto (f, e, -h)$ is an
automorphism of $\mathfrak{sl}_2$, so $e' = \Lambda$, $f' = L$, $h' = -H$ is again a representation, in which $\Lambda^k$ has weight $n - k$.
Corollary~\ref{representation-theory:corollary-lefschetz-structure} for $(e', f', h')$ says: $f'^{\,j} = L^j$ is an isomorphism from weight
$j$ to weight $-j$, which for $j = n - k$ is (1); the primitive vectors $\ker e' = \ker\Lambda$ occur only in weights $j \geq 0$, that is, in
degrees $k \leq n$; and
\[
\Lambda^\bullet = \bigoplus_{j \geq 0}\ \bigoplus_{i=0}^{j}L^iP^{n-j},
\]
which, restricted to degree $k$, is (2). For (3), write $\alpha = \sum_iL^i\alpha_i$ with $\alpha_i \in P^{k-2i}$. A primitive vector of
degree $m \leq n$ generates an irreducible representation with basis $\alpha, L\alpha, \ldots, L^{n-m}\alpha$ and $L^{n-m+1}\alpha = 0$
(Theorem~\ref{representation-theory:theorem-sl2-irreducibles}). So $L^{n-k+1}L^i\alpha_i = 0$ for $i = 0$, while for $i \geq 1$ the exponent
$n - k + 1 + i \leq n - (k - 2i)$ and $L^{n-k+1+i}\alpha_i \neq 0$ if $\alpha_i \neq 0$. These terms lie in different summands of the direct
sum above, so $L^{n-k+1}\alpha = 0$ if and only if $\alpha_i = 0$ for all $i \geq 1$, that is, if and only if $\alpha = \alpha_0$ is primitive.

For (4), $L$ has type $(1,1)$ because $\omega$ does (Proposition~\ref{kahler:proposition-fundamental-form-coordinates}); since the spaces
$\Lambda^{p,q}$ are mutually orthogonal (the basis $dz^I \wedge d\bar z^K$ of an orthonormal complex frame is orthogonal), the adjoint $\Lambda$ has
type $(-1,-1)$. For $\theta \in \RR$ let $T_\theta$ act on $\Lambda^{p,q}$ by $e^{i(p-q)\theta}$. It commutes with $L$ and $\Lambda$, hence maps
$L^jP^{k-2j}$ to itself, and by uniqueness in (2) $(T_\theta\alpha)_j = T_\theta(\alpha_j)$. If $\alpha$ has type $(p,q)$ then
$T_\theta\alpha = e^{i(p-q)\theta}\alpha$, so $T_\theta\alpha_j = e^{i(p-q)\theta}\alpha_j$ for all $\theta$, and $\alpha_j$ has only components of
types $(p', q')$ with $p' - q' = p - q$ and $p' + q' = k - 2j$, that is, $(p', q') = (p - j, q - j)$.
\end{proof}

\begin{remark}\label{kahler:remark-hodge-riemann-pointwise}
For a primitive form $\alpha \in P^k$ of type $(p,q)$, the Hodge star operator is given by \emph{Weil's formula}
\[
*\alpha = (-1)^{k(k+1)/2}\,i^{p-q}\,\frac{L^{n-k}\alpha}{(n-k)!},
\]
and it follows that $i^{p-q}(-1)^{k(k-1)/2}\,\alpha \wedge \bar\alpha \wedge \omega^{n-k}$ is a positive multiple of the volume form when
$\alpha \neq 0$ \cite[Section 1.2]{huybrechts-complex}. For instance, when $n = 1$ and $\alpha = dz$, one has $*dz = -i\,dz$ and
$i\,dz \wedge d\bar z = 2\,dx \wedge dy$. Integrated over a compact K\"ahler manifold, this positivity becomes the Hodge--Riemann bilinear relations of
Section~\ref{kahler:section-lefschetz-theorems}, the source of the polarisations of the Hodge structures of Part~\ref{part-hodge-theory}.
\end{remark}

\section{The Hodge Decomposition}\label{kahler:section-hodge-decomposition}
On a compact K\"ahler manifold the complex cohomology splits as
\[
H^k(X; \CC) = \bigoplus_{p+q=k}H^{p,q}(X), \qquad \overline{H^{p,q}(X)} = H^{q,p}(X), \qquad H^{p,q}(X) \cong H^q(X, \Omega^p_X)
\]
(Theorem~\ref{kahler:theorem-hodge-decomposition}). This \emph{Hodge decomposition} is the basic structure of the whole
subject: Part~\ref{part-hodge-theory} axiomatises it as a Hodge structure, and the Hodge conjecture asks which rational classes of type
$(p,p)$ come from algebraic cycles. The proof combines the Hodge theorem of Chapter~\ref{hodge-theorem} with the \emph{K\"ahler identities}
(Theorem~\ref{kahler:theorem-kahler-identities}), which give $\Delta_d = 2\Delta_{\bar\partial}$, so that the Laplacian
preserves types. We also prove the $\partial\bar\partial$-lemma, derive the topological restrictions on K\"ahler manifolds, and compute
examples. References: \cite[Chapter 0, Sections 6--7]{griffiths-harris}, \cite[Chapter 3]{huybrechts-complex},
\cite[Chapter 6]{voisin-hodge-1}, \cite[Chapter V]{wells-differential-analysis}.

\noindent\textbf{Prerequisites.} Chapter~\ref{kahler} (K\"ahler Manifolds) and Chapter~\ref{hodge-theorem} (Harmonic Forms and the Hodge
Theorem).

\section{Conventions}\label{kahler:section-conventions}

\begin{notation}\label{kahler:notation-hermitian-forms}
Let $X$ be a complex manifold of complex dimension $n$ with a Hermitian metric $g$ and fundamental form $\omega$. Complex forms carry the
Hermitian inner product $\langle\alpha, \beta\rangle$ extending the real one of Definition~\ref{riemannian-geometry:definition-hodge-star}
(linear in $\alpha$, conjugate linear in $\beta$). The Hodge star $*$ is extended complex linearly, so that
$\alpha \wedge *\bar\beta = \langle\alpha, \beta\rangle\,\mathrm{vol}_g$, and when $X$ is compact the $L^2$ product is
$(\alpha, \beta) = \int_X\alpha \wedge *\bar\beta$. The operators $d$, $\delta$ and $\Delta = \Delta_d = d\delta + \delta d$ are extended
complex linearly; since they are real, $\delta$ is still the formal adjoint of $d$. We write $\partial^*$, $\bar\partial^*$ for the formal
adjoints of $\partial$, $\bar\partial$, and
\[
\Delta_\partial = \partial\partial^* + \partial^*\partial, \qquad \Delta_{\bar\partial} = \bar\partial\bar\partial^* + \bar\partial^*\bar\partial,
\]
so that $\Delta_{\bar\partial}$ is the operator $\Box$ of Definition~\ref{complex-manifolds:definition-dbar-laplacian}. The operators $L$,
$\Lambda$, $H$ are those of Definition~\ref{kahler:definition-lefschetz-operators}, applied pointwise, and $\Pi^{p,q}$ is the projection
onto forms of type $(p,q)$. Finally $c(\alpha) = \bar\alpha$ is complex conjugation of forms. It is conjugate linear and satisfies
$(c\alpha, c\beta) = \overline{(\alpha, \beta)}$, $c\partial c = \bar\partial$, and $cLc = L$, $c\Lambda c = \Lambda$, since
$\omega$ is real.
\end{notation}

\begin{lemma}\label{kahler:lemma-adjoints-by-type}
The complex linear Hodge star maps $\Lambda^{p,q}$ to $\Lambda^{n-q,n-p}$, forms of different types are orthogonal, and
\[
\partial^* = -*\bar\partial*, \qquad \bar\partial^* = -*\partial*, \qquad \delta = \partial^* + \bar\partial^*,
\]
where $\partial^*$ has type $(-1, 0)$ and $\bar\partial^*$ has type $(0, -1)$. Moreover $c\,\partial^*c = \bar\partial^*$.
\end{lemma}

\begin{proof}
At a point choose holomorphic coordinates with $h_{j\bar k} = \delta_{jk}$ there (as in the proof of
Proposition~\ref{kahler:proposition-volume}). Then $dz^j = dx^j + i\,dy^j$ with $dx^j, dy^j$ orthonormal, so
$\langle dz^j, dz^k\rangle = 2\delta_{jk}$ and $\langle dz^j, d\bar z^k\rangle = \langle dx^j, dx^k\rangle - \langle dy^j, dy^k\rangle = 0$
(the term $i \cdot \overline{(-i)}$ is $-1$). Hence the products $dz^I \wedge d\bar z^J$ are pairwise orthogonal, and forms of different
types are orthogonal.

If $\beta \in \Lambda^{p,q}$, then $\alpha \wedge *\beta = \langle\alpha, \bar\beta\rangle\,\mathrm{vol} = 0$ for every $\alpha$ of
degree $p + q$ and type other than $(q,p)$, since $\bar\beta$ has type $(q,p)$. The wedge pairing
$\Lambda^{r,s} \times \Lambda^{n-r,n-s} \to \Lambda^{n,n}$ is perfect and pairs $\Lambda^{r,s}$ trivially with all other types of
complementary degree, so the component of $*\beta$ of type $(n-r, n-s)$ vanishes for every $(r,s) \neq (q,p)$. Thus $*\beta$ has type
$(n-q, n-p)$.

The real dimension of $X$ is even, so $\delta = -*d*$ by Definition~\ref{riemannian-geometry:definition-laplacian}. On a form of type
$(p,q)$, the operator $-*\partial*$ has type $(0, -1)$, since $(p,q) \mapsto (n-q, n-p) \mapsto (n-q+1, n-p) \mapsto (p, q-1)$, and
likewise $-*\bar\partial*$ has type $(-1, 0)$. So $\delta = D' + D''$ with $D' = -*\bar\partial*$ of type $(-1,0)$ and $D'' = -*\partial*$
of type $(0,-1)$. For $\alpha$ of type $(p-1, q)$ and $\beta$ of type $(p,q)$, orthogonality of types gives
$(\partial\alpha, \beta) = (d\alpha, \beta) = (\alpha, \delta\beta) = (\alpha, D'\beta)$, so $D' = \partial^*$; similarly $D'' = \bar\partial^*$.
Finally $(c\,\partial^*c\,\alpha, \beta) = \overline{(\partial^*c\alpha, c\beta)} = \overline{(c\alpha, \partial c\beta)} = (\alpha, c\,\partial c\,\beta) = (\alpha, \bar\partial\beta)$.
\end{proof}

\section{The K\"ahler identities}\label{kahler:section-kahler-identities}

\begin{lemma}\label{kahler:lemma-osculation-principle}
Let $P$ be a differential operator on forms, defined on every Hermitian manifold of dimension $n$, whose coefficients at each point, in
any holomorphic coordinates, are given by universal expressions in the values at that point of the $h_{j\bar k}$ and their first
derivatives. If $P = 0$ for the flat metric on $\CC^n$, then $P = 0$ on every K\"ahler manifold of dimension $n$.
\end{lemma}

\begin{proof}
Let $p \in X$ and choose holomorphic coordinates centred at $p$ with $h_{j\bar k} = \delta_{jk} + O(|z|^2)$
(Proposition~\ref{kahler:proposition-kahler-equivalent}). In this chart, the metric of $X$ and the flat metric $h_{j\bar k} = \delta_{jk}$ have
the same values and the same first derivatives at $p$. So the coefficients of $P$ at $p$ are the same for both metrics, and those of
the flat metric vanish.
\end{proof}

\begin{lemma}\label{kahler:lemma-car}
Let $W$ be a finite-dimensional complex vector space with a Hermitian inner product, and extend the inner product to $\Lambda^\bullet W$
as above. For $\theta \in W$ let $\varepsilon(\theta) = \theta \wedge -$ and let $\varepsilon(\theta)^*$ be its adjoint. Then
\[
\varepsilon(\theta)\varepsilon(\theta') = -\varepsilon(\theta')\varepsilon(\theta), \qquad
\varepsilon(\varphi)^*\varepsilon(\theta) + \varepsilon(\theta)\varepsilon(\varphi)^* = \langle\theta, \varphi\rangle\,\mathrm{id}.
\]
\end{lemma}

\begin{proof}
The first identity holds because $\theta \wedge \theta' = -\theta' \wedge \theta$. Both sides of the second are linear in $\theta$ and
conjugate linear in $\varphi$, so it suffices to take $\theta = e_a$, $\varphi = e_b$ from an orthonormal basis $e_1, \ldots, e_m$ and to
evaluate on $e_I = e_{i_1} \wedge \cdots \wedge e_{i_r}$, $i_1 < \cdots < i_r$. The adjoint $\varepsilon(e_b)^*$ sends $e_I$ to
$(-1)^{s-1}e_{I \setminus b}$ if $b = i_s \in I$, and to $0$ if $b \notin I$. If $a = b$, exactly one of the two terms is nonzero on
$e_I$, depending on whether $a \in I$, and it equals $e_I$. If $a \neq b$, both terms vanish unless $b \in I$ and $a \notin I$. In that case
moving $e_b$ into place in $e_a \wedge e_I$ costs one extra sign compared with $e_I$, so the two terms cancel.
\end{proof}

\begin{theorem}[K\"ahler identities]\label{kahler:theorem-kahler-identities}
On a K\"ahler manifold,
\[
[\Lambda, \bar\partial] = -i\partial^*, \qquad [\Lambda, \partial] = i\bar\partial^*, \qquad [\bar\partial^*, L] = i\partial, \qquad [\partial^*, L] = -i\bar\partial,
\]
and $[L, d] = 0$, $[\Lambda, \delta] = 0$.
\end{theorem}

\begin{proof}
The operator $P = [\Lambda, \bar\partial] + i\partial^*$ satisfies the hypothesis of
Lemma~\ref{kahler:lemma-osculation-principle}. Indeed $\Lambda$ and $*$ are algebraic operators whose coefficients are
universal functions of the $h_{j\bar k}$, and $P$ is obtained from them and from $\bar\partial$ by composition, using
$\partial^* = -*\bar\partial*$ (Lemma~\ref{kahler:lemma-adjoints-by-type}). The derivatives in $\bar\partial$ fall at most
once on these coefficients. So it suffices to prove $P = 0$ on $\CC^n$ with the flat metric. There $\omega = \frac{i}{2}\sum_kdz^k \wedge d\bar z^k$.
Write $e_k = \varepsilon(dz^k)$ and $\bar e_k = \varepsilon(d\bar z^k)$, which are constant operators. Then
\[
L = \frac{i}{2}\sum_ke_k\bar e_k, \qquad \Lambda = L^* = -\frac{i}{2}\sum_k\bar e_k^*e_k^*, \qquad \bar\partial = \sum_l\bar e_l\,\partial_{\bar z^l},
\qquad \partial = \sum_ke_k\,\partial_{z^k},
\]
where $\partial_{z^k}$, $\partial_{\bar z^l}$ act on the coefficients and commute with the constant operators. By
Lemma~\ref{kahler:lemma-car} and the inner products $\langle d\bar z^l, d\bar z^k\rangle = 2\delta_{kl}$,
$\langle d\bar z^l, dz^k\rangle = 0$ computed in Lemma~\ref{kahler:lemma-adjoints-by-type}, we have
$e_k^*\bar e_l = -\bar e_le_k^*$ and $\bar e_k^*\bar e_l = 2\delta_{kl} - \bar e_l\bar e_k^*$. Hence
\[
\bar e_k^*e_k^*\bar e_l = -\bar e_k^*\bar e_le_k^* = -2\delta_{kl}e_k^* + \bar e_l\bar e_k^*e_k^*,
\qquad\text{that is,}\qquad [\bar e_k^*e_k^*, \bar e_l] = -2\delta_{kl}e_k^*,
\]
and therefore $[\Lambda, \bar\partial] = -\frac{i}{2}\sum_k(-2)e_k^*\partial_{\bar z^k} = i\sum_ke_k^*\partial_{\bar z^k}$. On the other hand,
integration by parts against compactly supported forms shows that the formal adjoint of $\partial_{z^k}$ acting on coefficients is
$-\partial_{\bar z^k}$, because $\int\partial_{z}f\cdot\bar u = -\int f\,\overline{\partial_{\bar z}u}$. So
$\partial^* = \sum_k(\partial_{z^k})^*e_k^* = -\sum_ke_k^*\partial_{\bar z^k}$, and $[\Lambda, \bar\partial] = -i\partial^*$.

Conjugating with $c$ and using $c\Lambda c = \Lambda$, $c\bar\partial c = \partial$, $c\,\partial^*c = \bar\partial^*$ and $c(-i)c = i$ gives
$[\Lambda, \partial] = i\bar\partial^*$. Taking formal adjoints of the first two identities gives
$[\bar\partial^*, L] = (-i\partial^*)^* = i\partial$ and $[\partial^*, L] = (i\bar\partial^*)^* = -i\bar\partial$. Finally,
$d(\omega \wedge \alpha) = \omega \wedge d\alpha$ because $d\omega = 0$ and $\omega$ has even degree, so $[L, d] = 0$, and its adjoint is
$[\Lambda, \delta] = 0$.
\end{proof}

\begin{remark}\label{kahler:remark-identities-conventions}
The signs in the K\"ahler identities depend on conventions: the normalisation $\omega = \frac{i}{2}\sum h_{j\bar k}dz^j \wedge d\bar z^k$
used here, the sign of $\delta$, and whether inner products are linear in the first or in the second variable. They agree with
\cite[Proposition 3.1.12]{huybrechts-complex}. The proof via Lemma~\ref{kahler:lemma-osculation-principle} is the typical use
of the osculation property of K\"ahler metrics. The identities fail for non-K\"ahler Hermitian metrics, where torsion terms involving
$d\omega$ appear.
\end{remark}

\section{Comparison of Laplacians}\label{kahler:section-laplacians}

\begin{theorem}\label{kahler:theorem-laplacians}
On a K\"ahler manifold, $\partial\bar\partial^* + \bar\partial^*\partial = 0$, $\bar\partial\partial^* + \partial^*\bar\partial = 0$, and
\[
\Delta_d = 2\Delta_{\bar\partial} = 2\Delta_\partial .
\]
Consequently $\Delta_d$ preserves types: it commutes with every $\Pi^{p,q}$, with $\partial$, $\bar\partial$, $\partial^*$,
$\bar\partial^*$, and with $c$. Moreover $\Delta_d$ commutes with $*$, $L$ and $\Lambda$.
\end{theorem}

\begin{proof}
By Theorem~\ref{kahler:theorem-kahler-identities}, $i\bar\partial^* = [\Lambda, \partial]$, so
\[
i(\partial\bar\partial^* + \bar\partial^*\partial) = \partial\Lambda\partial - \partial\partial\Lambda + \Lambda\partial\partial - \partial\Lambda\partial = 0,
\]
using $\partial^2 = 0$. The second identity is its conjugate. Expanding,
\[
\Delta_d = (\partial + \bar\partial)(\partial^* + \bar\partial^*) + (\partial^* + \bar\partial^*)(\partial + \bar\partial)
= \Delta_\partial + \Delta_{\bar\partial} + (\partial\bar\partial^* + \bar\partial^*\partial) + (\bar\partial\partial^* + \partial^*\bar\partial) = \Delta_\partial + \Delta_{\bar\partial}.
\]
Using $\partial^* = i[\Lambda, \bar\partial]$ and $\bar\partial^* = -i[\Lambda, \partial]$,
\[
\Delta_\partial = i(\partial\Lambda\bar\partial - \partial\bar\partial\Lambda + \Lambda\bar\partial\partial - \bar\partial\Lambda\partial), \qquad
\Delta_{\bar\partial} = -i(\bar\partial\Lambda\partial - \bar\partial\partial\Lambda + \Lambda\partial\bar\partial - \partial\Lambda\bar\partial),
\]
so $\Delta_\partial - \Delta_{\bar\partial} = i\left(-(\partial\bar\partial + \bar\partial\partial)\Lambda + \Lambda(\bar\partial\partial + \partial\bar\partial)\right) = 0$ by
Proposition~\ref{complex-manifolds:proposition-dbar}. Hence $\Delta_d = 2\Delta_{\bar\partial} = 2\Delta_\partial$.

$\Delta_{\bar\partial}$ has type $(0,0)$, so $\Delta_d$ commutes with the $\Pi^{p,q}$. $\Delta_d$ commutes with $d$ and $\delta$, hence with
their components of each type, which are $\partial, \bar\partial, \partial^*, \bar\partial^*$. It is real, so it commutes with $c$, and it
commutes with $*$ by Proposition~\ref{riemannian-geometry:proposition-adjoint}. For $L$: by the K\"ahler identities
$[L, \delta] = -[\partial^*, L] - [\bar\partial^*, L] = i(\bar\partial - \partial)$ and $[L, d] = 0$, so
\[
[L, \Delta_d] = d[L, \delta] + [L, \delta]d = i\left((\partial + \bar\partial)(\bar\partial - \partial) + (\bar\partial - \partial)(\partial + \bar\partial)\right) = 0,
\]
again because $\partial^2 = \bar\partial^2 = \partial\bar\partial + \bar\partial\partial = 0$. Taking adjoints, $[\Lambda, \Delta_d] = 0$.
\end{proof}

\section{The Hodge decomposition}\label{kahler:section-hodge-decomposition-theorem}

\begin{definition}\label{kahler:definition-hpq}
For a complex manifold $X$, let $H^{p,q}(X) \subseteq H^{p+q}(X; \CC)$ be the subspace of classes that can be represented by a closed form
of type $(p,q)$. For a Hermitian metric, let $\mathcal{H}^k_{\CC}$ be the complex $\Delta_d$-harmonic $k$-forms and
$\mathcal{H}^{p,q} = \mathcal{H}^{p+q}_{\CC} \cap \mathcal{A}^{p,q}(X)$.
\end{definition}

\begin{theorem}[Hodge decomposition]\label{kahler:theorem-hodge-decomposition}
Let $X$ be a compact K\"ahler manifold. Then
\begin{enumerate}
\item $\mathcal{H}^k_{\CC} = \bigoplus_{p+q=k}\mathcal{H}^{p,q}$, and $\mathcal{H}^{p,q}$ is also the space of $\Delta_{\bar\partial}$-harmonic,
and of $\Delta_\partial$-harmonic, forms of type $(p,q)$;
\item $H^k(X; \CC) = \bigoplus_{p+q=k}H^{p,q}(X)$, and $\mathcal{H}^{p,q} \to H^{p,q}(X)$, $\alpha \mapsto [\alpha]$, is an isomorphism;
\item $\overline{H^{p,q}(X)} = H^{q,p}(X)$ (\emph{Hodge symmetry});
\item $H^{p,q}(X) \cong H^{p,q}_{\bar\partial}(X) \cong H^q(X, \Omega^p_X)$.
\end{enumerate}
The subspaces $H^{p,q}(X)$ depend only on the complex structure of $X$, not on the K\"ahler metric.
\end{theorem}

\begin{proof}
(1) If $\Delta_d\alpha = 0$ then $\Delta_d\Pi^{p,q}\alpha = \Pi^{p,q}\Delta_d\alpha = 0$ by
Theorem~\ref{kahler:theorem-laplacians}, so the components of a harmonic form are harmonic. The second statement holds
because $\Delta_d = 2\Delta_{\bar\partial} = 2\Delta_\partial$.

(2) Applying Theorem~\ref{hodge-theorem:theorem-hodge} to real and imaginary parts, and Theorem~\ref{de-rham:theorem-de-rham}, the map
$\mathcal{H}^k_{\CC} \to H^k(X; \CC)$ is an isomorphism. By (1), $H^k(X; \CC)$ is the direct sum of the images of the $\mathcal{H}^{p,q}$,
and each image lies in $H^{p,q}(X)$ because harmonic forms are closed. Conversely, let $\alpha$ be a closed form of type $(p,q)$ and let
$h$ be its harmonic projection. By Theorem~\ref{hodge-theorem:theorem-decomposition} (applied to real and imaginary parts),
$\alpha = h + d\beta$ for some $\beta$, and $h$ is the orthogonal projection of $\alpha$ onto $\mathcal{H}^k_{\CC}$. Since forms of different
types are orthogonal (Lemma~\ref{kahler:lemma-adjoints-by-type}), $\alpha$ is orthogonal to $\mathcal{H}^{r,s}$ for
$(r,s) \neq (p,q)$, so $h \in \mathcal{H}^{p,q}$ by (1). So $H^{p,q}(X)$ is exactly the image of $\mathcal{H}^{p,q}$, and (2) follows.

(3) $\Delta_d$ is a real operator, so $c$ maps $\mathcal{H}^{p,q}$ onto $\mathcal{H}^{q,p}$. Alternatively, the conjugate of a closed
$(p,q)$-form is a closed $(q,p)$-form.

(4) By (1), $\mathcal{H}^{p,q}$ is the space of $\Delta_{\bar\partial}$-harmonic $(p,q)$-forms, which is isomorphic to
$H^{p,q}_{\bar\partial}(X) \cong H^q(X, \Omega^p_X)$ by Theorem~\ref{complex-manifolds:theorem-dbar-hodge}.

The definition of $H^{p,q}(X)$ does not involve a metric.
\end{proof}

\begin{corollary}\label{kahler:corollary-hodge-numbers}
Let $X$ be a compact K\"ahler manifold of dimension $n$, with Hodge numbers $h^{p,q} = \dim H^q(X, \Omega^p_X)$. Then:
\begin{enumerate}
\item $b_k = \sum_{p+q=k}h^{p,q}$, so the Fr\"olicher spectral sequence degenerates at $E_1$;
\item $h^{p,q} = h^{q,p} = h^{n-p,n-q} = h^{n-q,n-p}$;
\item $b_{2k+1}$ is even for every $k$;
\item $h^{p,p} \geq 1$ for $0 \leq p \leq n$.
\end{enumerate}
\end{corollary}

\begin{proof}
(1) is Theorem~\ref{kahler:theorem-hodge-decomposition}(2) and (4) together with
Proposition~\ref{complex-manifolds:proposition-frolicher-inequality}. In (2), the first equality is Hodge symmetry and the second is Serre duality
(Theorem~\ref{complex-manifolds:theorem-serre-duality}). For (3), $b_{2k+1} = \sum_{p+q=2k+1}h^{p,q} = 2\sum_{p<q,\ p+q=2k+1}h^{p,q}$, since
$p \neq q$ when $p + q$ is odd. For (4), $\omega^p$ is a closed form of type $(p,p)$ and $[\omega]^p \neq 0$ by
Proposition~\ref{kahler:proposition-kahler-class-powers}.
\end{proof}

\begin{remark}\label{kahler:remark-hodge-diamond}
The Hodge numbers are displayed in the \emph{Hodge diamond}, with $h^{n,n}$ at the top, $h^{0,0}$ at the bottom, and $h^{p,q}$ in row
$p + q$. Row $k$ adds up to $b_k$, and the diamond is symmetric under the reflections $h^{p,q} \leftrightarrow h^{q,p}$ (Hodge symmetry) and
$h^{p,q} \leftrightarrow h^{n-p,n-q}$ (Serre duality). For a compact K\"ahler surface it reads
\[
\begin{array}{ccccc}
 & & h^{2,2} & & \\
 & h^{2,1} & & h^{1,2} & \\
h^{2,0} & & h^{1,1} & & h^{0,2} \\
 & h^{1,0} & & h^{0,1} & \\
 & & h^{0,0} & &
\end{array}
\qquad = \qquad
\begin{array}{ccccc}
 & & 1 & & \\
 & q & & q & \\
p_g & & h^{1,1} & & p_g \\
 & q & & q & \\
 & & 1 & &
\end{array}
\]
where $q = h^{0,1}$ is the \emph{irregularity} and $p_g = h^{2,0}$ the \emph{geometric genus}.
\end{remark}

\begin{remark}\label{kahler:remark-hodge-filtration}
The \emph{Hodge filtration} $F^pH^k(X; \CC) = \bigoplus_{r \geq p}H^{r,k-r}(X)$ determines the decomposition, since
$H^{p,q} = F^p \cap \overline{F^q}$ when $p + q = k$. Because the Fr\"olicher spectral sequence degenerates, $F^p$ is also the image of the
hypercohomology of the truncated complex $\Omega^{\geq p}_X$ in $H^k(X; \CC)$ \cite[Chapter 8]{voisin-hodge-1}. The filtration varies
holomorphically in families, whereas the decomposition does not. This is why the theory of Hodge structures in
Part~\ref{part-hodge-theory} is phrased with filtrations.
\end{remark}

\begin{proposition}\label{kahler:proposition-holomorphic-forms}
On a compact K\"ahler manifold every holomorphic $p$-form is closed and harmonic, and
$H^0(X, \Omega^p_X) = \mathcal{H}^{p,0} \cong H^{p,0}(X) \subseteq H^p(X; \CC)$. In particular a holomorphic form that is exact, or
merely cohomologous to zero, vanishes, and $b_1(X) = 2\dim H^0(X, \Omega^1_X)$.
\end{proposition}

\begin{proof}
If $\alpha$ is a holomorphic $p$-form, then $\bar\partial\alpha = 0$ and $\bar\partial^*\alpha = 0$, the latter because
$\bar\partial^*\alpha$ would have type $(p, -1)$. So $\Delta_{\bar\partial}\alpha = 0$, hence $\Delta_d\alpha = 0$ and $\alpha$ is closed
and harmonic. Conversely, if $\alpha \in \mathcal{H}^{p,0}$ then
$0 = (\Delta_{\bar\partial}\alpha, \alpha) = \|\bar\partial\alpha\|^2 + \|\bar\partial^*\alpha\|^2$, so $\bar\partial\alpha = 0$ and $\alpha$ is holomorphic. The
harmonic representative of a class is unique, so a harmonic form cohomologous to zero vanishes. Finally $b_1 = h^{1,0} + h^{0,1} = 2h^{1,0}$.
\end{proof}

\begin{proposition}\label{kahler:proposition-functoriality}
Let $f \colon X \to Y$ be a holomorphic map of compact K\"ahler manifolds. Then $f^*H^{p,q}(Y) \subseteq H^{p,q}(X)$. The cup product
maps $H^{p,q} \times H^{r,s}$ into $H^{p+r,q+s}$. If $\dim X = n$, the Poincar\'e duality pairing
$H^k(X; \CC) \times H^{2n-k}(X; \CC) \to \CC$ restricts to perfect pairings $H^{p,q}(X) \times H^{n-p,n-q}(X) \to \CC$, and it is zero on
$H^{p,q} \times H^{r,s}$ for all other pairs of types.
\end{proposition}

\begin{proof}
Pullback by a holomorphic map and wedge product preserve types of forms, and the wedge product corresponds to the cup product
(Theorem~\ref{de-rham:theorem-de-rham}). For $(r,s) \neq (n-p, n-q)$ with $r + s = 2n - p - q$, the wedge product of forms of types
$(p,q)$ and $(r,s)$ has type $(p+r, q+s) \neq (n,n)$ and so is zero, being a $2n$-form. The pairing on $H^k \times H^{2n-k}$ is perfect
(Theorem~\ref{de-rham:theorem-poincare-duality-de-rham}), and by
Theorem~\ref{kahler:theorem-hodge-decomposition} it is block diagonal with blocks $H^{p,q} \times H^{n-p,n-q}$. Each block
is therefore perfect.
\end{proof}

\begin{proposition}\label{kahler:proposition-cycle-classes-type}
Let $X$ be a compact K\"ahler manifold of dimension $n$ and $Z \subseteq X$ an analytic subvariety of pure codimension $p$. Then its cycle
class $[Z] \in H^{2p}(X; \CC)$ (Proposition~\ref{complex-manifolds:proposition-subvariety-class}) lies in $H^{p,p}(X)$.
\end{proposition}

\begin{proof}
Let $k = n - p = \dim Z$. The restriction of a form of type $(r,s)$ to the complex manifold $Z_{\mathrm{reg}}$ of dimension $k$ has type
$(r,s)$ there, so it vanishes if $r + s = 2k$ and $(r,s) \neq (k,k)$. Hence $\int_Z\beta = 0$ for every closed form $\beta$ of type
$(r,s) \neq (k,k)$ with $r + s = 2k$. By definition of the cycle class, $\int_X\eta \wedge \beta = \int_Z\beta$ for a closed form $\eta$
representing $[Z]$ (Theorem~\ref{currents:theorem-currents}). Write $[Z] = \sum_{a+b=2p}z^{a,b}$ with $z^{a,b} \in H^{a,b}(X)$. By
Proposition~\ref{kahler:proposition-functoriality}, $z^{a,b}$ pairs only with $H^{n-a,n-b}$, and it pairs to zero with it
unless $(n-a, n-b) = (k,k)$. So $z^{a,b}$ pairs to zero with everything unless $(a,b) = (p,p)$, and by perfectness $z^{a,b} = 0$ for
$(a,b) \neq (p,p)$.
\end{proof}

\begin{remark}\label{kahler:remark-hodge-classes}
The cycle class of $Z$ is the image of an integral class in $H^{2p}(X; \ZZ)$ (Chapter~\ref{intersection-theory}). So algebraic cycles give
\emph{Hodge classes}, the elements of $H^{2p}(X; \QQ) \cap H^{p,p}(X)$. The Hodge conjecture (Chapter~\ref{hodge-conjecture}) predicts that
on a projective manifold every Hodge class is a rational linear combination of cycle classes.
\end{remark}

\section{The $\partial\bar\partial$-lemma}\label{kahler:section-ddbar-lemma}

\begin{lemma}\label{kahler:lemma-del-hodge}
Let $X$ be compact K\"ahler. Then
$\mathcal{A}^{p,q}(X) = \mathcal{H}^{p,q} \oplus \partial\mathcal{A}^{p-1,q}(X) \oplus \partial^*\mathcal{A}^{p+1,q}(X)$ orthogonally, and
$\ker\partial = \mathcal{H}^{p,q} \oplus \partial\mathcal{A}^{p-1,q}(X)$. The same holds with $\bar\partial$, $\bar\partial^*$ in place of
$\partial$, $\partial^*$.
\end{lemma}

\begin{proof}
For $\bar\partial$ the decomposition is Theorem~\ref{complex-manifolds:theorem-dbar-hodge}, with $\mathcal{H}^{p,q}$ the $\Delta_{\bar\partial}$-harmonic
forms (Theorem~\ref{kahler:theorem-hodge-decomposition}(1)). Applying the conjugate linear isometry $c$ to the decomposition
of $\mathcal{A}^{q,p}(X)$, and using $c\bar\partial c = \partial$, $c\bar\partial^*c = \partial^*$ and $c\mathcal{H}^{q,p} = \mathcal{H}^{p,q}$, gives
the decomposition for $\partial$. If $\partial\alpha = 0$ and $\alpha = h + \partial\beta + \partial^*\gamma$, then
$0 = (\partial\alpha, \gamma) = (\partial\partial^*\gamma, \gamma) = \|\partial^*\gamma\|^2$, because $\partial h = 0$ and $\partial^2 = 0$. So
$\alpha \in \mathcal{H}^{p,q} \oplus \operatorname{im}\partial$; the argument for $\bar\partial$ is the same.
\end{proof}

\begin{theorem}[$\partial\bar\partial$-lemma]\label{kahler:theorem-ddbar-lemma}
Let $X$ be a compact K\"ahler manifold and $\alpha$ a $d$-closed form of type $(p,q)$. The following are equivalent:
\begin{enumerate}
\item $\alpha$ is $d$-exact;
\item $\alpha$ is $\partial$-exact;
\item $\alpha$ is $\bar\partial$-exact;
\item $\alpha = \partial\bar\partial\gamma$ for a form $\gamma$ of type $(p-1, q-1)$.
\end{enumerate}
\end{theorem}

\begin{proof}
(4) implies the others: $\partial\bar\partial\gamma = d(\bar\partial\gamma) = \partial(\bar\partial\gamma) = \bar\partial(-\partial\gamma)$.

Assume (1), (2) or (3). Then $\alpha$ is orthogonal to $\mathcal{H}^{p,q}$: for $h$ harmonic, $(d\beta, h) = (\beta, \delta h) = 0$, and
likewise with $\partial^*h = 0$ or $\bar\partial^*h = 0$ (harmonic forms are killed by all of these, by
Theorem~\ref{kahler:theorem-hodge-decomposition}(1)). Since $\partial\alpha = 0$ (the $(p+1,q)$-part of $d\alpha = 0$),
Lemma~\ref{kahler:lemma-del-hodge} gives $\alpha = \partial\gamma$ with $\gamma$ of type $(p-1,q)$. Decompose
$\gamma = h' + \bar\partial\mu + \bar\partial^*\nu$ by the same lemma. Then $\partial h' = 0$, and
$\partial\bar\partial^*\nu = -\bar\partial^*\partial\nu$ by Theorem~\ref{kahler:theorem-laplacians}, so
\[
\alpha = \partial\bar\partial\mu - \bar\partial^*\partial\nu .
\]
Now $\bar\partial\alpha = 0$ and $\bar\partial\partial\bar\partial\mu = -\partial\bar\partial\bar\partial\mu = 0$, hence
$\bar\partial\bar\partial^*\partial\nu = 0$ and $\|\bar\partial^*\partial\nu\|^2 = (\bar\partial\bar\partial^*\partial\nu, \partial\nu) = 0$. Therefore
$\alpha = \partial\bar\partial\mu$, with $\mu$ of type $(p-1, q-1)$.
\end{proof}

\begin{corollary}\label{kahler:corollary-canonical-dolbeault}
Let $X$ be compact K\"ahler. Every Dolbeault class in $H^{p,q}_{\bar\partial}(X)$ has a $d$-closed representative, and
$[\alpha] \mapsto [\alpha]_{\bar\partial}$, for $\alpha$ a closed $(p,q)$-form, is a well-defined isomorphism
$H^{p,q}(X) \to H^{p,q}_{\bar\partial}(X)$ that does not depend on any metric.
\end{corollary}

\begin{proof}
If $\alpha$ and $\alpha'$ are closed $(p,q)$-forms in the same de Rham class, then $\alpha - \alpha'$ is a $d$-exact closed form of type
$(p,q)$, hence $\bar\partial$-exact by Theorem~\ref{kahler:theorem-ddbar-lemma}. So the map is well defined. For a K\"ahler
metric, restricting it to harmonic representatives gives the isomorphism
$\mathcal{H}^{p,q} \cong H^{p,q}_{\bar\partial}(X)$ of Theorem~\ref{complex-manifolds:theorem-dbar-hodge}. The harmonic representative of a Dolbeault
class is $d$-closed.
\end{proof}

\begin{corollary}\label{kahler:corollary-ddbar-real}
Let $X$ be compact K\"ahler and let $\alpha$, $\alpha'$ be closed real $(1,1)$-forms with $[\alpha] = [\alpha']$ in $H^2(X; \RR)$. Then
$\alpha' = \alpha + i\partial\bar\partial\varphi$ for a smooth real function $\varphi$. In particular any two K\"ahler forms with the same
K\"ahler class differ by $i\partial\bar\partial\varphi$.
\end{corollary}

\begin{proof}
By Theorem~\ref{kahler:theorem-ddbar-lemma}, $\alpha' - \alpha = \partial\bar\partial\mu$ for a complex function $\mu$. Since
$\alpha' - \alpha$ is real, $\partial\bar\partial\mu = \overline{\partial\bar\partial\mu} = \bar\partial\partial\bar\mu = -\partial\bar\partial\bar\mu$. So
$\alpha' - \alpha = \frac12\partial\bar\partial(\mu - \bar\mu) = i\partial\bar\partial\operatorname{Im}\mu$, and $\varphi = \operatorname{Im}\mu$ works.
\end{proof}

\begin{remark}\label{kahler:remark-formality}
Deligne, Griffiths, Morgan and Sullivan deduced from the $\partial\bar\partial$-lemma that the real homotopy type of a compact K\"ahler
manifold is determined by its real cohomology ring (the manifold is \emph{formal}). In particular all Massey products vanish, which gives
obstructions to K\"ahlerness beyond the Hodge numbers \cite{dgms}. The $\partial\bar\partial$-lemma, and hence the Hodge decomposition, also
hold for compact complex manifolds bimeromorphic to K\"ahler ones \cite{dgms}.
\end{remark}

\section{Examples}\label{kahler:section-examples-hodge-decomposition}

\begin{example}\label{kahler:example-curves}
Let $X$ be a compact Riemann surface. It is K\"ahler (Example~\ref{kahler:example-kahler}), so $b_1$ is even, and its Hodge diamond has $h^{0,0} = h^{1,1} = 1$ and
$h^{1,0} = h^{0,1} = g$, where $g = b_1/2$. So the \emph{genus} $g$, a topological invariant, equals both the dimension of the space of
holomorphic $1$-forms $H^0(X, K_X)$ and $\dim H^1(X, \mathcal{O}_X)$. For the Riemann sphere $g = 0$; for an elliptic curve
$\CC/\Lambda$, $g = 1$ and $H^0(X, K_X)$ is spanned by $dz$.
\end{example}

\begin{example}\label{kahler:example-projective-space}
For $\CC\PP^n$, $b_{2k} = 1$ for $0 \leq k \leq n$ and all other Betti numbers vanish
(Theorem~\ref{singular-homology:theorem-projective-spaces}). By Corollary~\ref{kahler:corollary-hodge-numbers}(4),
$h^{k,k} \geq 1$, so $H^{2k}(\CC\PP^n; \CC) = H^{k,k}$ is spanned by $[\omega_{FS}]^k$, and $h^{p,q} = 0$ for $p \neq q$. In particular
$\CC\PP^n$ has no nonzero holomorphic forms of positive degree, and $H^q(\CC\PP^n, \mathcal{O}) = 0$ for $q > 0$.
\end{example}

\begin{example}\label{kahler:example-elliptic-periods}
Let $E_\tau = \CC/(\ZZ + \tau\ZZ)$ with $\operatorname{Im}\tau > 0$, and let $a$, $b$ be the classes in $H_1(E_\tau; \ZZ)$ of the images of
the segments $[0, 1]$ and $[0, \tau]$. Integration identifies $H^1(E_\tau; \CC)$ with $\CC^2$ through
$[\alpha] \mapsto (\int_a\alpha, \int_b\alpha)$. Then $H^{1,0}$ is spanned by $[dz] \mapsto (1, \tau)$ and $H^{0,1}$ by
$[d\bar z] \mapsto (1, \bar\tau)$. The lattice $H^1(E_\tau; \ZZ)$ and the vector space $H^1(E_\tau; \CC)$ do not depend on $\tau$, but the
line $H^{1,0} = [1 : \tau] \in \PP^1$ does, and it recovers $\tau$. So the Hodge decomposition captures the complex structure, which is not a
topological invariant. This is the first example of a \emph{period map} (Chapter~\ref{vhs}).
\end{example}

\begin{example}\label{kahler:example-quartic}
For a smooth quartic surface $X \subseteq \CC\PP^3$ (a K3 surface) the Hodge diamond is
\[
\begin{array}{ccccc}
 & & 1 & & \\
 & 0 & & 0 & \\
1 & & 20 & & 1 \\
 & 0 & & 0 & \\
 & & 1 & &
\end{array}
\]
Here $b_1 = 0$ by the Lefschetz hyperplane theorem (Chapter~\ref{kodaira}), $K_X \cong \mathcal{O}_X$ by adjunction
(Proposition~\ref{complex-manifolds:proposition-adjunction}) so that $p_g = 1$, and $b_2 = 22$ comes from the Euler characteristic
$\chi = c_2(X) = 24$ \cite[Chapter VIII]{bhpv}. The subspace $H^{1,1} \cap H^2(X; \QQ)$ is where the Hodge conjecture for
divisors, the Lefschetz $(1,1)$-theorem, lives.
\end{example}

\begin{example}\label{kahler:example-product-curves}
Let $C$ and $C'$ be compact Riemann surfaces of genera $g$ and $g'$, so $h^{0,0} = h^{1,1} = 1$ and $h^{1,0} = h^{0,1} = g$ for $C$ (Example~\ref{kahler:example-curves}). By
Exercise~\ref{kahler:exercise-kunneth-hodge}, the surface $S = C \times C'$ has
\[
h^{1,0}(S) = g + g', \qquad h^{2,0}(S) = gg', \qquad h^{1,1}(S) = 2 + 2gg',
\]
since $H^{1,1}(S)$ receives $H^{1,1}(C) \otimes H^{0,0}(C')$, $H^{0,0}(C) \otimes H^{1,1}(C')$, $H^{1,0}(C) \otimes H^{0,1}(C')$ and $H^{0,1}(C) \otimes H^{1,0}(C')$. As a check, $b_2(S) = 2h^{2,0} + h^{1,1} = 2 + 4gg'$, which is
$1 + (2g)(2g') + 1$ by the K\"unneth formula. For two elliptic curves, $g = g' = 1$, the Hodge diamond is
\[
\begin{array}{ccccc}
 & & 1 & & \\
 & 2 & & 2 & \\
1 & & 4 & & 1 \\
 & 2 & & 2 & \\
 & & 1 & &
\end{array}
\]
which is also that of every two-dimensional complex torus (Example~\ref{complex-manifolds:example-hodge-numbers}). The $H^{1,1}$ of $S$ always contains the classes of the two fibres, and the remaining $2gg'$
dimensions come from $H^1(C) \otimes H^1(C')$; which rational classes lie there is governed by the homomorphisms between the Jacobians of $C$ and $C'$ (for elliptic curves this is worked out in
Chapter~\ref{hodge-conjecture}).
\end{example}

\begin{counterexample}\label{kahler:counterexample-hopf-hodge-decomposition}
The Hopf surface (Exercise~\ref{complex-manifolds:exercise-hopf-surface}) has $b_1 = 1$, which is odd, so by
Corollary~\ref{kahler:corollary-hodge-numbers}(3) it is not K\"ahler, which gives a second proof of
Counterexample~\ref{kahler:counterexample-hopf}. It also shows how the theorem fails: $h^{1,0} = 0 \neq 1 = h^{0,1}$
(Example~\ref{complex-manifolds:example-hodge-numbers}).
\end{counterexample}

\begin{counterexample}\label{kahler:counterexample-iwasawa}
Let $G$ be the complex Heisenberg group of matrices
$\left(\begin{smallmatrix}1 & z_1 & z_3 \\ 0 & 1 & z_2 \\ 0 & 0 & 1\end{smallmatrix}\right)$, $z \in \CC^3$, and $\Gamma \subseteq G$ the
subgroup with entries in $\ZZ[i]$. The \emph{Iwasawa manifold} $X = \Gamma\backslash G$ is a compact complex manifold of dimension $3$ carrying
a holomorphic $1$-form that is not closed. So $X$ is not K\"ahler, by Proposition~\ref{kahler:proposition-holomorphic-forms}.
\end{counterexample}

\begin{proof}
Left multiplication by an element with entries $(a_1, a_2, a_3)$ is
$(z_1, z_2, z_3) \mapsto (z_1 + a_1, z_2 + a_2, z_3 + a_1z_2 + a_3)$. This map is holomorphic, and $\Gamma$ acts freely and properly
discontinuously, being a discrete subgroup acting by translations of the group. So $X$ is a complex manifold. It is compact: using
$a_1, a_2 \in \ZZ[i]$ one moves $z_1, z_2$ into the unit square, and then using $a_3$ one moves $z_3$ there, so the image of a compact
set surjects onto $X$. The form $\theta = dz_3 - z_1\,dz_2$ is invariant under left multiplication:
$d(z_3 + a_1z_2 + a_3) - (z_1 + a_1)\,d(z_2 + a_2) = dz_3 - z_1\,dz_2$. So it descends to a holomorphic $1$-form on $X$, and
$d\theta = -dz_1 \wedge dz_2 \neq 0$.
\end{proof}

\section{Hard Lefschetz and the Hodge--Riemann Relations}\label{kahler:section-lefschetz-theorems}
Let $X$ be a compact K\"ahler manifold of dimension $n$ with K\"ahler form $\omega$. Cup product with the K\"ahler class is an operator
$L$ on $H^\bullet(X; \RR)$ of degree $2$. The \emph{hard Lefschetz theorem}
(Theorem~\ref{kahler:theorem-hard-lefschetz}) says that
\[
L^{n-k} \colon H^k(X; \RR) \to H^{2n-k}(X; \RR)
\]
is an isomorphism for $k \leq n$. The proof takes the pointwise $\mathfrak{sl}_2$-action of Theorem~\ref{kahler:theorem-sl2-forms} and
passes it to harmonic forms, using that $L$ and $\Lambda$ commute with the Laplacian. The result is a canonical decomposition of cohomology
into primitive pieces. On them the intersection form, twisted by powers of $\omega$, has definite signs: these are the
\emph{Hodge--Riemann bilinear relations} (Theorem~\ref{kahler:theorem-hodge-riemann}). Their consequences include the
\emph{Hodge index theorem} (Theorem~\ref{kahler:theorem-hodge-index}) and, in Part~\ref{part-hodge-theory}, the polarisations of
the Hodge structures of projective manifolds. References: \cite[Chapter 0, Section 7]{griffiths-harris}, \cite[Section 3.3]{huybrechts-complex},
\cite[Chapter 6]{voisin-hodge-1}.

\noindent\textbf{Prerequisites.} Section~\ref{kahler:section-hodge-decomposition} (The Hodge Decomposition).

Throughout, $X$ is a compact K\"ahler manifold of dimension $n$ with K\"ahler form $\omega$, and we use the notation of
Notation~\ref{kahler:notation-hermitian-forms}. We identify $H^k(X; \CC)$ with the space $\mathcal{H}^k_\CC$ of complex
harmonic forms and $H^k(X; \RR)$ with the real harmonic forms (Theorem~\ref{hodge-theorem:theorem-hodge}).

\section{The $\mathfrak{sl}_2$-action on cohomology}\label{kahler:section-sl2-cohomology}

\begin{definition}\label{kahler:definition-lefschetz-cohomology}
The \emph{Lefschetz operator} on cohomology is $L \colon H^k(X; \RR) \to H^{k+2}(X; \RR)$, $L(a) = [\omega] \cup a$, which is
represented on forms by $\alpha \mapsto \omega \wedge \alpha$. A class $a \in H^k(X; \RR)$ with $k \leq n$ is \emph{primitive} if
$L^{n-k+1}a = 0$. We write $P^k(X) = P^k(X; \RR)$ for the primitive classes of degree $k$, $P^k(X; \CC)$ for their complexification,
and $P^{p,q}(X) = P^{p+q}(X; \CC) \cap H^{p,q}(X)$.
\end{definition}

\begin{theorem}[Hard Lefschetz]\label{kahler:theorem-hard-lefschetz}
Let $X$ be a compact K\"ahler manifold of dimension $n$ with K\"ahler class $[\omega]$. Then:
\begin{enumerate}
\item the pointwise operators $L$, $\Lambda$ and $H$ preserve the space $\mathcal{H}^\bullet$ of harmonic forms, which becomes a
finite-dimensional representation of $\mathfrak{sl}_2$ in which $\mathcal{H}^k$ has weight $k - n$;
\item for $0 \leq k \leq n$, $L^{n-k} \colon H^k(X; \RR) \to H^{2n-k}(X; \RR)$ is an isomorphism;
\item (\emph{Lefschetz decomposition}) for every $k$,
\[
H^k(X; \RR) = \bigoplus_{j \geq 0}L^jP^{k-2j}(X),
\]
and $L^j$ is injective on $P^{k-2j}(X)$;
\item a class of degree $k \leq n$ is primitive if and only if its harmonic representative $\alpha$ satisfies $\Lambda\alpha = 0$, that is,
$\alpha$ is primitive at every point;
\item the decompositions are compatible with the Hodge decomposition:
$P^k(X; \CC) = \bigoplus_{p+q=k}P^{p,q}(X)$ and $H^{p,q}(X) = \bigoplus_jL^jP^{p-j,q-j}(X)$.
\end{enumerate}
\end{theorem}

\begin{proof}
(1) By Theorem~\ref{kahler:theorem-laplacians}, $L$ and $\Lambda$ commute with $\Delta_d$, so they preserve
$\mathcal{H}^\bullet = \ker\Delta_d$; so does $H$, which acts on $\mathcal{H}^k$ as multiplication by $k - n$. The commutation relations
$[L, \Lambda] = H$, $[H, L] = 2L$, $[H, \Lambda] = -2\Lambda$ hold pointwise by Theorem~\ref{kahler:theorem-sl2-forms}, hence on
$\mathcal{H}^\bullet$. The space $\mathcal{H}^\bullet$ is finite-dimensional by Theorem~\ref{hodge-theorem:theorem-decomposition}, and it
is real since all three operators are real.

(2) and (3) now follow exactly as statements (1) and (2) of Theorem~\ref{kahler:theorem-sl2-forms} were derived from
Corollary~\ref{representation-theory:corollary-lefschetz-structure}: that argument used only a finite-dimensional
$\mathfrak{sl}_2$-representation in which the degree $k$ part has weight $k - n$. Under the Hodge isomorphism, $L$ on harmonic forms
is the cup product with $[\omega]$, since $\omega \wedge \alpha$ is harmonic and represents $[\omega] \cup [\alpha]$
(Theorem~\ref{de-rham:theorem-de-rham}). This gives (2), and it gives (3) with the primitive subspaces defined as
$\ker\Lambda \cap \mathcal{H}^{k-2j}$.

(4) Part (3) of Theorem~\ref{kahler:theorem-sl2-forms} was also derived from the representation theory alone. Applied to
$\mathcal{H}^\bullet$, it says that for $\alpha \in \mathcal{H}^k$ with $k \leq n$, $\Lambda\alpha = 0$ if and only if
$L^{n-k+1}\alpha = 0$. So the subspaces $\ker\Lambda \cap \mathcal{H}^k$ used in (3) are the $P^k(X)$ of
Definition~\ref{kahler:definition-lefschetz-cohomology}.

(5) $L$ maps $H^{p,q}$ to $H^{p+1,q+1}$. So the kernel of $L^{n-k+1}$ on $H^k(X; \CC) = \bigoplus_{p+q=k}H^{p,q}$ is the direct sum of
its kernels on the summands, which gives the first statement. The second follows by comparing types in (3).
\end{proof}

\begin{remark}\label{kahler:remark-metric-independence}
Primitivity is defined by $L$ alone, so $P^k(X)$ and the Lefschetz decomposition depend only on the K\"ahler class $[\omega]$, not on the
metric. The operator $\Lambda$ on cohomology was defined using harmonic forms, but it too depends only on $[\omega]$, by
Proposition~\ref{kahler:proposition-lambda-unique}.
\end{remark}

\begin{proposition}\label{kahler:proposition-lambda-unique}
Let $V$ be a finite-dimensional vector space over a field $K$ of characteristic zero, graded by degree. Let $e$ and $h$ be endomorphisms
such that $h$ acts on the degree $k$ part by $k - n$ and $[h, e] = 2e$. Then there is at most one endomorphism $f$ with $[e, f] = h$ and
$[h, f] = -2f$. If $K \subseteq K'$ is a field extension and such an $f$ exists over $K'$, then it is defined over $K$. In particular, if $X$ is
projective and $[\omega]$ is the restriction of the hyperplane class, then $L$ and $\Lambda$ preserve $H^\bullet(X; \QQ)$, and the
Lefschetz decomposition is defined over $\QQ$.
\end{proposition}

\begin{proof}
It suffices to treat $K = \CC$ for uniqueness, after extending scalars. If $f$ and $f'$ both work, then $g = f - f'$ satisfies $[e, g] = 0$ and
$[h, g] = -2g$. The Lie algebra $\mathfrak{sl}_2$ acts on $\End(V)$ through the triple $(e, f, h)$ by $x \cdot \phi = [x, \phi]$. In this
representation $g$ is killed by $e$ and has weight $-2$. By Theorem~\ref{representation-theory:theorem-sl2-irreducibles}, a nonzero vector
killed by $e$ has a non-negative weight, so $g = 0$. The conditions on $f$ are linear equations with coefficients in $K$. A system of linear
equations over $K$ that has a solution over $K'$ has one over $K$, since the rank of a matrix does not change under field extension, and by
uniqueness the two solutions coincide. For the last statement, $[\omega] = c_1(\mathcal{O}(1))|_X$ lies in the image of
$H^2(X; \ZZ)$ (Proposition~\ref{kahler:proposition-fubini-study}), so $L$ preserves $H^\bullet(X; \QQ)$, and the rest follows.
\end{proof}

\begin{corollary}\label{kahler:corollary-betti-inequalities}
Let $X$ be a compact K\"ahler manifold of dimension $n$. Then $L \colon H^k \to H^{k+2}$ is injective for $k < n$ and surjective for
$k \geq n - 1$. Consequently
\[
b_k \leq b_{k+2} \quad (k < n), \qquad h^{p,q} \leq h^{p+1,q+1} \quad (p + q < n), \qquad \dim P^k = b_k - b_{k-2} \quad (k \leq n).
\]
\end{corollary}

\begin{proof}
For $k < n$, $L^{n-k}$ is injective on $H^k$ and equals $L^{n-k-1} \circ L$, so $L$ is injective on $H^k$. For $k \geq n - 1$, put
$j = 2n - k - 2 \leq n - 1$. Then $L^{n-j} = L \circ L^{n-j-1}$ maps $H^j$ onto $H^{2n-j} = H^{k+2}$, and $L^{n-j-1}$ lands in $H^k$, so $L$ is
onto $H^{k+2}$. The inequality for Hodge numbers follows since $L$ has type $(1,1)$, and the last formula from
$H^k = P^k \oplus LH^{k-2}$ with $L$ injective on $H^{k-2}$.
\end{proof}

\begin{example}\label{kahler:example-tori}
Let $X = \CC^n/\Lambda$ be a complex torus with a flat K\"ahler metric, so that $\omega$ has constant coefficients. The harmonic forms are the
forms with constant coefficients (Example~\ref{complex-manifolds:example-hodge-numbers}), so $H^\bullet(X; \RR) \cong \Lambda^\bullet V^*$ with
$V = \RR^{2n}$, and $L$ is the wedge product with the constant form $\omega$. In this case
Theorem~\ref{kahler:theorem-hard-lefschetz} is exactly the linear algebra of Theorem~\ref{kahler:theorem-sl2-forms}. For
instance, for $n = 2$, $H^2(X; \RR)$ has dimension $6$ and is the sum of $\RR[\omega]$ and the $5$-dimensional primitive part
(Exercise~\ref{kahler:exercise-primitive-dimension}).
\end{example}

\begin{example}\label{kahler:example-betti-obstruction}
The manifold $S^2 \times S^3 \times S^3$ has Poincar\'e polynomial $(1 + t^2)(1 + t^3)^2 = 1 + t^2 + 2t^3 + 2t^5 + t^6 + t^8$. Its real
dimension is $8$, so a K\"ahler structure would have $n = 4$, and Corollary~\ref{kahler:corollary-betti-inequalities} would
give $b_2 \leq b_4$. But $b_2 = 1 > 0 = b_4$, so no complex structure on this manifold is K\"ahler.
\end{example}

\section{Pointwise Hodge--Riemann relations}\label{kahler:section-pointwise}

\begin{proposition}\label{kahler:proposition-pointwise-hodge-riemann}
Let $V$ be a complex vector space of dimension $n$ with a Hermitian inner product and fundamental form $\omega$, as in
Definition~\ref{kahler:definition-lefschetz-operators}. Let $\alpha$ be a primitive form of type $(p,q)$ and degree $k = p + q \leq n$. Then
\[
i^{p-q}(-1)^{k(k-1)/2}\,\alpha \wedge \bar\alpha \wedge \omega^{n-k} = (n-k)!\,|\alpha|^2\,\mathrm{vol} .
\]
More generally, if $\alpha$ and $\beta$ are primitive of type $(p,q)$, then
$i^{p-q}(-1)^{k(k-1)/2}\alpha \wedge \bar\beta \wedge \omega^{n-k} = (n-k)!\,\langle\alpha, \beta\rangle\,\mathrm{vol}$.
\end{proposition}

\begin{proof}
The second statement follows from the first by polarisation, since both sides are Hermitian forms in $(\alpha, \beta)$. Choose a basis
with $h_{j\bar k} = \delta_{jk}$, so that $dz_j = dx_j + i\,dy_j$ with $dx_j, dy_j$ orthonormal, and put
$\nu_j = dx_j \wedge dy_j = \frac{i}{2}dz_j \wedge d\bar z_j$. Then $\omega = \sum_j\nu_j$, and $\mathrm{vol} = \nu_1 \wedge \cdots \wedge \nu_n$.
For a set $R$ of indices write $\nu_R = \bigwedge_{j \in R}\nu_j$, and for disjoint $A, B$ write $D = A \cup B$ and
$dz_A \wedge d\bar z_B$ for the wedge products in increasing order.

\emph{Step 1: splitting.} Every monomial $dz_I \wedge d\bar z_J$ equals $\pm(2/i)^{|I \cap J|}dz_A \wedge d\bar z_B \wedge \nu_C$ with
$C = I \cap J$, $A = I \setminus C$, $B = J \setminus C$. So $\Lambda^{p,q}$ is the orthogonal direct sum of the subspaces
$dz_A \wedge d\bar z_B \wedge S_D$, where $S_D$ is spanned by the $\nu_R$ with $R \cap D = \emptyset$. Write
$\alpha = \sum_{A,B}dz_A \wedge d\bar z_B \wedge f_{AB}$ with $f_{AB} \in S_D$ homogeneous of degree $r = p - |A| = q - |B|$ in the $\nu_j$. In the proof
of Theorem~\ref{kahler:theorem-sl2-forms}, $\Lambda = \sum_j\Lambda_j$, where $\Lambda_j$ acts on the $j$-th factor, kills $1$, $dz_j$
and $d\bar z_j$, and sends $\nu_j$ to $1$. Hence
$\Lambda(dz_A \wedge d\bar z_B \wedge f) = dz_A \wedge d\bar z_B \wedge \Lambda f$, where $\Lambda\nu_R = \sum_{j \in R}\nu_{R \setminus j}$, and
$\alpha$ is primitive if and only if every $f_{AB}$ is.

\emph{Step 2: the summands are orthogonal for the left-hand side.} Let $\alpha$ and $\beta$ lie in the summands of $(A, B)$ and
$(A', B')$. For $j \in A$, the form $\alpha$ contains $dz_j$ but not $d\bar z_j$. A top-degree monomial needs exactly one $d\bar z_j$, and it
cannot come from a power of $\omega$ or from a $\nu_j$ in $\bar\beta$, both of which would repeat $dz_j$. So it must come from
$\bar\beta = d\bar z_{A'} \wedge dz_{B'} \wedge \bar f'$, which forces $j \in A'$. Thus $\alpha \wedge \bar\beta \wedge \omega^{n-k} = 0$
unless $A \subseteq A'$, and by symmetry $A = A'$ and $B = B'$. The summands are also orthogonal for $\langle\ ,\ \rangle$, so it suffices
to treat $\alpha = dz_A \wedge d\bar z_B \wedge f$ with $|A| = a$, $|B| = b$, $f = \sum_{|R|=r}f_R\nu_R \in S_D$ primitive, and $m = n - a - b$
indices outside $D$. Then $n - k = m - 2r$.

\emph{Step 3: the $\nu$-part.} Only the $\nu_j$ with $j \notin D$ survive in $\omega^{n-k}$ after wedging with $dz_A \wedge d\bar z_B \wedge d\bar z_A \wedge dz_B$.
Moreover $\big(\sum_{j \notin D}\nu_j\big)^{m-2r} = (m-2r)!\sum_{|T|=m-2r}\nu_T$, and $\nu_R\nu_{R'}\nu_T$ is the product of all $\nu_j$, $j \notin D$, exactly
when $R, R', T$ partition the complement of $D$. So
\[
f \wedge \bar f \wedge \Big(\sum_{j \notin D}\nu_j\Big)^{m-2r} = (m-2r)!\sum_{R \cap R' = \emptyset}f_R\overline{f_{R'}}\;\nu_{\{1, \ldots, n\} \setminus D}.
\]
We claim that $\sum_{R \cap R' = \emptyset}f_R\overline{f_{R'}} = (-1)^r\sum_R|f_R|^2$ for primitive $f$. Primitivity says that
$\sum_{R' \supseteq Q}f_{R'} = 0$ for every set $Q$ of size $r - 1$. For $|S| = s < r$, each $R' \supseteq S$ of size $r$ contains exactly $r - s$ sets
$Q$ of size $r - 1$ with $S \subseteq Q$, so $\sum_{R' \supseteq S}f_{R'} = \frac{1}{r-s}\sum_Q\sum_{R' \supseteq Q}f_{R'} = 0$. By
inclusion--exclusion, $[R \cap R' = \emptyset] = \sum_{S \subseteq R \cap R'}(-1)^{|S|}$. Hence
\[
\sum_{R'}\overline{f_{R'}}\,[R \cap R' = \emptyset] = \sum_{S \subseteq R}(-1)^{|S|}\sum_{R' \supseteq S}\overline{f_{R'}} = (-1)^r\,\overline{f_R},
\]
because only $S = R$ survives. Multiplying by $f_R$ and summing over $R$ proves the claim.

\emph{Step 4: the $dz$-part and the signs.} Reordering gives
\[
dz_A \wedge d\bar z_B \wedge d\bar z_A \wedge dz_B = (-1)^{ab + b}(dz_A \wedge d\bar z_A) \wedge (dz_B \wedge d\bar z_B)
= (-1)^{ab + b + a(a-1)/2 + b(b-1)/2}(-2i)^{a+b}\,\nu_D,
\]
using $dz_A \wedge d\bar z_A = (-1)^{a(a-1)/2}\bigwedge_{j \in A}(dz_j \wedge d\bar z_j)$ and $dz_j \wedge d\bar z_j = -2i\,\nu_j$. The forms $f$ and $\bar f$
have even degree, so they can be moved freely. Therefore $i^{p-q}(-1)^{k(k-1)/2}\alpha \wedge \bar\alpha \wedge \omega^{n-k}$ equals
$(m-2r)!\sum_R|f_R|^2\,\mathrm{vol}$ times
\[
i^{a-b}(-2i)^{a+b}(-1)^{k(k-1)/2 + ab + b + a(a-1)/2 + b(b-1)/2 + r} .
\]
Now $i^{a-b}(-i)^{a+b} = (-1)^{a+b}i^{2a} = (-1)^b$. With $s = a + b$ and $k = s + 2r$, we have
$k(k-1)/2 \equiv s(s-1)/2 + r \equiv a(a-1)/2 + b(b-1)/2 + ab + r \pmod 2$. So the total sign is $+1$ and the factor is $2^{a+b}$. Finally
$|\alpha|^2 = |dz_A \wedge d\bar z_B|^2|f|^2 = 2^{a+b}\sum_R|f_R|^2$, since the $dz_j$, $d\bar z_j$ are orthogonal of squared length $2$
(Lemma~\ref{kahler:lemma-adjoints-by-type}) and the $\nu_R$ are orthonormal. As $m - 2r = n - k$, this is the assertion.
\end{proof}

\begin{remark}\label{kahler:remark-weil-formula}
Proposition~\ref{kahler:proposition-pointwise-hodge-riemann} also follows from Weil's formula
$*\alpha = (-1)^{k(k+1)/2}i^{p-q}L^{n-k}\alpha/(n-k)!$ for primitive $\alpha$ of type $(p,q)$ (Remark~\ref{kahler:remark-hodge-riemann-pointwise}),
using $\beta \wedge *\bar\alpha = \langle\beta, \alpha\rangle\,\mathrm{vol}$. The direct computation above avoids proving Weil's formula in
general.
\end{remark}

\section{The Hodge--Riemann bilinear relations}\label{kahler:section-hodge-riemann}

\begin{definition}\label{kahler:definition-q}
For $k \leq n$ define the bilinear form $Q$ on $H^k(X; \RR)$, and its complex bilinear extension to $H^k(X; \CC)$, by
\[
Q(a, b) = (-1)^{k(k-1)/2}\int_Xa \cup b \cup [\omega]^{n-k} .
\]
\end{definition}

\begin{theorem}[Hodge--Riemann bilinear relations]\label{kahler:theorem-hodge-riemann}
Let $X$ be a compact K\"ahler manifold of dimension $n$ and $k \leq n$. Then:
\begin{enumerate}
\item $Q(b, a) = (-1)^kQ(a, b)$, and $Q$ is nondegenerate on $H^k(X; \RR)$;
\item $Q(H^{p,q}, H^{p',q'}) = 0$ unless $(p', q') = (q, p)$;
\item the Lefschetz decomposition is $Q$-orthogonal: $Q(L^jP^{k-2j}, L^{j'}P^{k-2j'}) = 0$ for $j \neq j'$;
\item for $0 \neq a \in P^{p,q}(X)$, $p + q = k$, we have $i^{p-q}Q(a, \bar a) > 0$.
\end{enumerate}
\end{theorem}

\begin{proof}
(1) The symmetry is graded commutativity of the cup product. For nondegeneracy, if $Q(a, b) = 0$ for all $b$, then $a \cup [\omega]^{n-k}$
pairs trivially with $H^k$, so $L^{n-k}a = 0$ by Poincar\'e duality (Theorem~\ref{de-rham:theorem-poincare-duality-de-rham}), and $a = 0$ by
Theorem~\ref{kahler:theorem-hard-lefschetz}.

(2) is Proposition~\ref{kahler:proposition-functoriality}, as $L^{n-k}$ has type $(n-k, n-k)$.

(3) Let $j < j'$, $a \in P^{k-2j}$ and $b \in P^{k-2j'}$. Then $Q(L^ja, L^{j'}b) = \pm\int_Xb \cup L^{n-k+j+j'}a$, and
$n - k + j + j' \geq n - (k - 2j) + 1$, so $L^{n-k+j+j'}a = 0$ by primitivity. The case $j > j'$ follows by (1).

(4) Let $\alpha$ be the harmonic representative of $a$. It has type $(p,q)$ (Theorem~\ref{kahler:theorem-hodge-decomposition})
and is primitive at every point (Theorem~\ref{kahler:theorem-hard-lefschetz}(4)), and $\bar\alpha$ represents $\bar a$. By
Proposition~\ref{kahler:proposition-pointwise-hodge-riemann},
\[
i^{p-q}Q(a, \bar a) = \int_Xi^{p-q}(-1)^{k(k-1)/2}\alpha \wedge \bar\alpha \wedge \omega^{n-k} = (n-k)!\,\|\alpha\|^2 > 0 .
\]
\end{proof}

\begin{remark}\label{kahler:remark-polarisation}
In the language of Chapter~\ref{hodge-structures}: if $[\omega]$ is integral, then $P^k(X; \QQ)$ with the decomposition
$P^k(X; \CC) = \bigoplus P^{p,q}$ and the form $Q$ is a \emph{polarised Hodge structure} of weight $k$. By
Theorem~\ref{kahler:theorem-hodge-riemann}(3), $H^k(X; \QQ)$ is an orthogonal direct sum of such structures, with $Q$ twisted by
signs on the summands. This is why the Hodge structures of projective manifolds are polarisable, a fact that underlies the semisimplicity
arguments of Part~\ref{part-hodge-theory}.
\end{remark}

\begin{example}\label{kahler:example-curves-riemann}
For a compact Riemann surface ($n = 1$) and $k = 1$, every class is primitive, and part (4) says that $i\int_X\alpha \wedge \bar\alpha > 0$ for
every nonzero holomorphic $1$-form $\alpha$. These are \emph{Riemann's bilinear relations}. For $E_\tau = \CC/(\ZZ + \tau\ZZ)$ and $\alpha = dz$,
\[
i\int_{E_\tau}dz \wedge d\bar z = i\cdot(-2i)\operatorname{area}(E_\tau) = 2\operatorname{Im}\tau > 0,
\]
so the positivity condition $\operatorname{Im}\tau > 0$ on the period $\tau$ is an instance of the Hodge--Riemann relations. It is the
condition cutting out the period domain, here the upper half plane (Chapter~\ref{vhs}).
\end{example}

\section{The Hodge index theorem}\label{kahler:section-hodge-index}

\begin{theorem}[Hodge index theorem]\label{kahler:theorem-hodge-index}
Let $X$ be a compact K\"ahler manifold of even dimension $n$. The signature of the intersection form $(a, b) \mapsto \int_Xa \cup b$ on
$H^n(X; \RR)$ is
\[
\sigma(X) = \sum_{p,q = 0}^n(-1)^ph^{p,q}(X).
\]
\end{theorem}

\begin{proof}
Let $I(a, b) = \int_Xa \cup b$ on $H^n$. By Theorem~\ref{kahler:theorem-hodge-riemann}(3) and the Hodge decomposition, $H^n(X; \RR)$
is the $I$-orthogonal direct sum of the real subspaces $W_{j,p,q} = L^j\big(P^{p,q} + P^{q,p}\big) \cap H^n(X; \RR)$ over $j \geq 0$ and
$p \leq q$ with $p + q = n - 2j =: k$. Here $P^{p,q}$ is paired by $I$ only with $P^{q,p}$ (Theorem~\ref{kahler:theorem-hodge-riemann}(2)).
We compute the sign of $I$ on each piece. Note that $k$ is even. For $a \in P^{p,q}$ we have
$I(L^ja, L^j\bar a) = \int a \cup \bar a \cup [\omega]^{2j} = (-1)^{k(k-1)/2}Q(a, \bar a)$, since $n - k = 2j$. So
$I(L^ja, L^j\bar a) = (-1)^{k(k-1)/2}i^{q-p}\cdot i^{p-q}Q(a, \bar a)$ has the sign $\varepsilon_{p,q} = (-1)^{k(k-1)/2 + (q-p)/2}$. Writing $k = 2m$,
$k(k-1)/2 \equiv m$ and $(q - p)/2 = m - p$, so $\varepsilon_{p,q} = (-1)^p$.
\begin{itemize}
\item If $p < q$, the real points of $L^j(P^{p,q} \oplus P^{q,p})$ are the $L^j(a + \bar a)$, $a \in P^{p,q}$, and
$I(L^j(a + \bar a), L^j(a + \bar a)) = 2I(L^ja, L^j\bar a)$, because $I(L^ja, L^ja) = 0$ by type. So $I$ is definite of sign $(-1)^p$ on a
real space of dimension $2\dim P^{p,q}$.
\item If $p = q$, then $I$ is definite of sign $(-1)^p$ on the real points of $L^jP^{p,p}$, which has real dimension $\dim P^{p,p}$.
\end{itemize}
Hence $\sigma(X) = \sum_{p+q \leq n,\ p+q \text{ even}}(-1)^p\dim P^{p,q}$. By
Theorem~\ref{kahler:theorem-hard-lefschetz}(5) and Corollary~\ref{kahler:corollary-betti-inequalities},
$\dim P^{p,q} = h^{p,q} - h^{p-1,q-1}$ for $p + q \leq n$. Substituting $p' = p - 1$, $q' = q - 1$ in the second sum gives
\[
\sigma(X) = \sum_{p+q = n}(-1)^ph^{p,q} + 2\sum_{p+q < n,\ p+q \text{ even}}(-1)^ph^{p,q} .
\]
On the other hand, in $\sum_{p,q}(-1)^ph^{p,q}$ the terms with $p + q$ odd cancel in pairs by Hodge symmetry $h^{p,q} = h^{q,p}$, since then
$(-1)^p = -(-1)^q$. The terms with $p + q > n$ even equal those with $p + q < n$ by Serre duality $h^{p,q} = h^{n-p,n-q}$, since $n$ is
even. So the two expressions agree.
\end{proof}

\begin{corollary}\label{kahler:corollary-surfaces}
Let $X$ be a compact K\"ahler surface with irregularity $q = h^{0,1}$ and geometric genus $p_g = h^{2,0}$. The intersection form on
$H^2(X; \RR)$ has signature $(2p_g + 1, h^{1,1} - 1)$, and its restriction to $H^{1,1}_\RR = H^{1,1}(X) \cap H^2(X; \RR)$ has signature
$(1, h^{1,1} - 1)$. So $b^+(X) = 2p_g + 1$, and a real $(1,1)$-class $a$ with $\int_Xa \cup [\omega] = 0$ satisfies $\int_Xa^2 \leq 0$, with
equality only for $a = 0$.
\end{corollary}

\begin{proof}
For $n = 2$ the pieces in the proof of Theorem~\ref{kahler:theorem-hodge-index} are: $L(P^{0,0}) = \RR[\omega]$ with sign $+$;
$P^{0,2} + P^{2,0} = H^{0,2} \oplus H^{2,0}$ with sign $+$ and real dimension $2p_g$ (every class of type $(2,0)$ is primitive since
$H^{3,1} = 0$); and $P^{1,1}$ with sign $-$ and dimension $h^{1,1} - 1$. The real $(1,1)$-classes are $\RR[\omega] \oplus P^{1,1}_\RR$. A class
$a \in H^{1,1}_\RR$ with $\int a \cup [\omega] = 0$ is orthogonal to $[\omega]$, so it lies in $P^{1,1}_\RR$, where the form is negative
definite.
\end{proof}

\begin{remark}\label{kahler:remark-hodge-index-divisors}
For a projective surface, applied to classes of divisors, Corollary~\ref{kahler:corollary-surfaces} is the classical Hodge
index theorem: the intersection form on the N\'eron--Severi group has exactly one positive eigenvalue (Chapter~\ref{intersection-theory}).
Grothendieck's standard conjecture of Hodge type asks for the analogue of Theorem~\ref{kahler:theorem-hodge-riemann}(4) for
algebraic cycles in arbitrary characteristic (Chapter~\ref{absolute-hodge}).
\end{remark}

\section{Further remarks and counterexamples}\label{kahler:section-remarks}

\begin{remark}\label{kahler:remark-symplectic}
Hard Lefschetz is a special property of K\"ahler manifolds. Thurston gave a compact symplectic $4$-manifold, a nilmanifold with
$b_1 = 3$, that admits no K\"ahler structure \cite{thurston-1976}. Its first Betti number is odd, so
Corollary~\ref{kahler:corollary-hodge-numbers} already rules out a K\"ahler structure. Moreover, for every symplectic form
$\omega$ the Lefschetz map $[\omega] \cup - \colon H^1 \to H^3$ fails to be an isomorphism: otherwise $(a, b) \mapsto \int a \cup b \cup [\omega]$
would be a nondegenerate skew-symmetric form on the odd-dimensional space $H^1$. For singular projective varieties, hard Lefschetz
fails for ordinary cohomology but holds for intersection cohomology \cite{bbd}. In characteristic $p$, Deligne proved hard Lefschetz for
$\ell$-adic cohomology as part of his second proof of the Weil conjectures \cite{deligne-weil-2}.
\end{remark}

\begin{counterexample}\label{kahler:counterexample-non-primitive-sign}
The sign in Theorem~\ref{kahler:theorem-hodge-riemann}(4) genuinely requires primitivity. On $X = \CC\PP^1 \times \CC\PP^1$ with
$\omega = \mathrm{pr}_1^*\omega_{FS} + \mathrm{pr}_2^*\omega_{FS}$, let $h_1$, $h_2$ be the pullbacks of the generators, so $[\omega] = h_1 + h_2$,
$h_1^2 = h_2^2 = 0$ and $\int h_1h_2 = 1$. For $k = 2$, $Q(a, a) = -\int a^2$. Then $Q(h_1 - h_2, h_1 - h_2) = 2 > 0$ for the primitive class
$h_1 - h_2$, as predicted for type $(1,1)$, while $Q([\omega], [\omega]) = -2 < 0$ for the non-primitive class $[\omega]$.
\end{counterexample}

\section{Exercises}\label{kahler:section-exercises}

\subsection{K\"ahler Manifolds}\label{kahler:subsection-exercises-kahler-manifolds}

\begin{exercise}\label{kahler:exercise-fubini-study-line}
Compute $\int_{\CC\PP^1}\omega_{FS}$ directly in the affine coordinate $w$, and deduce that $\int_{\CC\PP^n}\omega_{FS}^n = 1$.
\end{exercise}

\begin{solution}
On $\CC\PP^1$, $\omega_{FS} = \frac{i}{2\pi}\partial\bar\partial\log(1 + |w|^2) = \frac{i}{2\pi}\frac{dw \wedge d\bar w}{(1 + |w|^2)^2}$, and the
point at infinity has measure zero. In polar coordinates $\frac{i}{2}dw \wedge d\bar w = r\,dr\,d\theta$, so
\[
\int_{\CC\PP^1}\omega_{FS} = \frac1\pi\int_0^{2\pi}\int_0^\infty\frac{r\,dr}{(1 + r^2)^2}\,d\theta = \frac1\pi \cdot 2\pi \cdot \frac12 = 1 .
\]
A line generates $H_2(\CC\PP^n; \ZZ)$ (it is the $2$-skeleton of the standard cell structure), so $[\omega_{FS}]$ is the generator $\alpha$
of $H^2(\CC\PP^n; \ZZ)$ that is positive on lines, and
$\alpha^n$ generates $H^{2n}(\CC\PP^n; \ZZ)$ (Theorem~\ref{singular-homology:theorem-projective-spaces}); it is positive on the
fundamental class since $\omega_{FS}^n$ is a positive multiple of the volume form (Proposition~\ref{kahler:proposition-volume}). So
$\int_{\CC\PP^n}\omega_{FS}^n = 1$.
\end{solution}

\begin{exercise}\label{kahler:exercise-primitive-dimension}
Show that $\dim P^k = \binom{2n}{k} - \binom{2n}{k-2}$ for $k \leq n$, and list the dimensions of $P^0, P^1, P^2$ when $n = 2$.
\end{exercise}

\begin{solution}
By Theorem~\ref{kahler:theorem-sl2-forms}(2), $\Lambda^k = P^k \oplus L\Lambda^{k-2}$, and $L \colon \Lambda^{k-2} \to \Lambda^k$ is injective
for $k \leq n$ by (1) (as $L^{n-k+2}$ is injective on $\Lambda^{k-2}$). For $n = 2$ the dimensions are $1$, $4$ and $6 - 1 = 5$.
\end{solution}

\begin{exercise}\label{kahler:exercise-no-global-potential}
Show that the K\"ahler form of a compact K\"ahler manifold of positive dimension is never of the form $i\partial\bar\partial\varphi$ for a
smooth function $\varphi$ on all of $X$, while on $\CC^n$ one has $\omega_0 = i\partial\bar\partial(\frac12|z|^2)$.
\end{exercise}

\begin{solution}
$i\partial\bar\partial\varphi = d(i\bar\partial\varphi)$, because $d = \partial + \bar\partial$ and $\bar\partial\bar\partial\varphi = 0$. So
$[\omega] = 0$, contradicting Proposition~\ref{kahler:proposition-kahler-class-powers} for $k = 1$. On $\CC^n$,
$i\partial\bar\partial\frac12\sum_jz^j\bar z^j = \frac{i}{2}\sum_jdz^j \wedge d\bar z^j = \omega_0$.
\end{solution}

\begin{exercise}\label{kahler:exercise-hopf-direct}
Let $C \subseteq X$ be the image of the line $\{z_1 = 0\}$ in the Hopf surface of Counterexample~\ref{kahler:counterexample-hopf}. Show that $C$ is
an elliptic curve and give a second proof that $X$ is not K\"ahler using $\int_C\omega$.
\end{exercise}

\begin{solution}
$C = (\CC^\times)/\langle 2\rangle \cong \CC/(2\pi i\ZZ + \log 2\,\ZZ)$ via the exponential, a complex torus of dimension one. If $\omega$ were a
K\"ahler form, then $\omega = d\eta$ because $H^2(X; \RR) = 0$, so $\int_C\omega = \int_Cd\eta = 0$ by Stokes' theorem; but
$\int_C\omega = \mathrm{vol}(C) > 0$ by Proposition~\ref{kahler:proposition-volume}.
\end{solution}

\begin{exercise}\label{kahler:exercise-kahler-cone}
Show that the set of K\"ahler classes of a compact K\"ahler manifold is an open convex cone in the real vector space
$H^{1,1}(X) \cap H^2(X; \RR)$ of classes represented by closed real $(1,1)$-forms.
\end{exercise}

\begin{solution}
If $\omega_1, \omega_2$ are K\"ahler forms and $s, t > 0$, then $s\omega_1 + t\omega_2$ is a closed positive real $(1,1)$-form, hence K\"ahler; this
gives convexity and the cone property. For openness, choose closed real $(1,1)$-forms $\beta_1, \ldots, \beta_m$ whose classes form a basis of the
space. Positivity is an open condition on the coefficient matrix and $X$ is compact, so $\omega + \sum_ic_i\beta_i$ is positive for all small
$|c_i|$.
\end{solution}

\begin{exercise}\label{kahler:exercise-product-lefschetz}
Let $V = V_1 \oplus V_2$ be an orthogonal sum of Hermitian vector spaces of complex dimensions $n_1$, $n_2$. Show that
$L = L_1 \otimes 1 + 1 \otimes L_2$, $\Lambda = \Lambda_1 \otimes 1 + 1 \otimes \Lambda_2$ on $\Lambda^\bullet V^* = \Lambda^\bullet V_1^* \otimes \Lambda^\bullet V_2^*$,
and that the $\mathfrak{sl}_2$-representation of Theorem~\ref{kahler:theorem-sl2-forms} is the tensor product of those for $V_1$ and $V_2$.
Deduce the decomposition of $\Lambda^\bullet(\CC^n)^*$ into irreducibles using the Clebsch--Gordan formula.
\end{exercise}

\begin{solution}
$\omega = \omega_1 + \omega_2$, and the tensor product inner product is the induced one, which gives the formulas for $L$ and $\Lambda$ (the
operators have even degree, so there are no signs). The counting operators satisfy $H = H_1 \otimes 1 + 1 \otimes H_2$ since $n = n_1 + n_2$;
this is the tensor product representation. For $n = 1$ the representation on $\Lambda^\bullet(\CC)^* \cong \CC^4$ is $V(1) \oplus 2V(0)$ (from
the proof of Theorem~\ref{kahler:theorem-sl2-forms}: $1, x \wedge y$ span $V(1)$, and $x$, $y$ are primitive of weight $0$). Hence
$\Lambda^\bullet(\CC^n)^* \cong (V(1) \oplus 2V(0))^{\otimes n}$, which decomposes by the Clebsch--Gordan formula $V(a) \otimes V(b) \cong V(a + b) \oplus V(a + b - 2) \oplus \cdots \oplus V(|a - b|)$
(Exercise~\ref{representation-theory:exercise-clebsch-gordan}).
\end{solution}

\subsection{The Hodge Decomposition}\label{kahler:subsection-exercises-hodge-decomposition}

\begin{exercise}\label{kahler:exercise-kunneth-hodge}
Let $X$ and $Y$ be compact K\"ahler manifolds. Show that
$h^{p,q}(X \times Y) = \sum h^{p_1,q_1}(X)\,h^{p_2,q_2}(Y)$, the sum over $p_1 + p_2 = p$, $q_1 + q_2 = q$.
\end{exercise}

\begin{solution}
$X \times Y$ is K\"ahler (Example~\ref{kahler:example-kahler}). The K\"unneth isomorphism
$H^\bullet(X; \CC) \otimes H^\bullet(Y; \CC) \to H^\bullet(X \times Y; \CC)$, $a \otimes b \mapsto \mathrm{pr}_X^*a \cup \mathrm{pr}_Y^*b$
(Theorem~\ref{singular-homology:theorem-kunneth}), maps $H^{p_1,q_1}(X) \otimes H^{p_2,q_2}(Y)$ into $H^{p_1+p_2,q_1+q_2}(X \times Y)$ by
Proposition~\ref{kahler:proposition-functoriality}. The source is the direct sum of these pieces and the target is the direct
sum of its $H^{p,q}$. So the isomorphism maps $\bigoplus_{p_1+p_2=p,\ q_1+q_2=q}H^{p_1,q_1} \otimes H^{p_2,q_2}$ onto $H^{p,q}(X \times Y)$.
\end{solution}

\begin{exercise}\label{kahler:exercise-calabi-eckmann}
Calabi and Eckmann constructed complex structures on $S^{2a+1} \times S^{2b+1}$ for all $a, b \geq 0$ \cite{calabi-eckmann}. Show that none
of them is K\"ahler, except on $S^1 \times S^1$.
\end{exercise}

\begin{solution}
The complex dimension is $a + b + 1 \geq 1$. By the K\"unneth formula, the cohomology of $S^{2a+1} \times S^{2b+1}$ lives in degrees
$0$, $2a + 1$, $2b + 1$ and $2a + 2b + 2$. If $a, b \geq 1$, then $H^2 = 0$, contradicting
Proposition~\ref{kahler:proposition-kahler-class-powers}. If $a = 0 < b$, then $b_1 = 1$ is odd, contradicting
Corollary~\ref{kahler:corollary-hodge-numbers}(3), and the case $b = 0 < a$ is symmetric. The remaining case $a = b = 0$ is a
torus of real dimension $2$, and every complex structure on it is K\"ahler (Example~\ref{kahler:example-kahler}).
\end{solution}

\begin{exercise}\label{kahler:exercise-invariant-fubini-study}
Show that $\omega_{FS}$ is the only $U(n+1)$-invariant K\"ahler form on $\CC\PP^n$ in its cohomology class.
\end{exercise}

\begin{solution}
Let $\omega'$ be another one. By Corollary~\ref{kahler:corollary-ddbar-real}, $\omega' = \omega_{FS} + i\partial\bar\partial\varphi$
with $\varphi$ real. For $g \in U(n+1)$, $g^*\omega' = \omega'$ and $g^*\omega_{FS} = \omega_{FS}$, so $i\partial\bar\partial(\varphi \circ g) = i\partial\bar\partial\varphi$,
since $g$ is holomorphic. Averaging over $U(n+1)$ with its normalised Haar measure gives an invariant function
$\tilde\varphi = \int\varphi \circ g\,dg$ with $i\partial\bar\partial\tilde\varphi = i\partial\bar\partial\varphi$ (differentiate under the integral). As
$U(n+1)$ acts transitively, $\tilde\varphi$ is constant. Hence $i\partial\bar\partial\varphi = 0$ and $\omega' = \omega_{FS}$.
\end{solution}

\begin{exercise}\label{kahler:exercise-h1-real}
Let $X$ be compact K\"ahler. Show that the map $H^1(X; \RR) \to H^1(X, \mathcal{O}_X)$ induced by the inclusion $\RR_X \to \mathcal{O}_X$ is an
isomorphism of real vector spaces. (This is the starting point for the Picard variety $H^1(X, \mathcal{O}_X)/H^1(X; \ZZ)$.)
\end{exercise}

\begin{solution}
Projecting forms to their $(0,q)$-components defines a morphism from the de Rham resolution of $\CC_X$ to the Dolbeault resolution of
$\mathcal{O}_X$ over the inclusion $\CC_X \to \mathcal{O}_X$, because the $(0, q+1)$-part of $d\alpha$ is $\bar\partial$ of the
$(0,q)$-part of $\alpha$. So the map sends $[\alpha]$ to $[\alpha^{0,1}]_{\bar\partial}$. Both sides have real dimension $b_1 = 2h^{0,1}$.
If $\alpha$ is a real harmonic $1$-form with $\alpha^{0,1}$ $\bar\partial$-exact, then $\alpha^{0,1}$ is harmonic (by
Theorem~\ref{kahler:theorem-hodge-decomposition}(1)) and $\bar\partial$-exact, hence zero. Then $\alpha^{1,0} = \overline{\alpha^{0,1}} = 0$,
so $\alpha = 0$ and the map is injective.
\end{solution}

\begin{exercise}\label{kahler:exercise-lefschetz-harmonic}
Let $X$ be compact K\"ahler. Show that $L$ and $\Lambda$ map harmonic forms to harmonic forms, and that the primitive components
(Theorem~\ref{kahler:theorem-sl2-forms}) of a harmonic form are harmonic.
\end{exercise}

\begin{solution}
$L$ and $\Lambda$ commute with $\Delta_d$ (Theorem~\ref{kahler:theorem-laplacians}), so they preserve $\ker\Delta_d$. Let
$\alpha = \sum_jL^j\alpha_j$ be the primitive decomposition of a harmonic form. Then $\Lambda\Delta_d\alpha_j = \Delta_d\Lambda\alpha_j = 0$,
so $\Delta_d\alpha_j$ is primitive, and $0 = \Delta_d\alpha = \sum_jL^j(\Delta_d\alpha_j)$ is a primitive decomposition of $0$. By
uniqueness, every $\Delta_d\alpha_j$ vanishes.
\end{solution}

\begin{exercise}\label{kahler:exercise-holomorphic-closed-surface}
Show that on a compact complex surface every holomorphic $1$-form is closed, even without assuming the surface is K\"ahler.
(Hint: for a holomorphic $1$-form $\theta$, compute $\int_Xd\theta \wedge \overline{d\theta}$ in two ways.)
\end{exercise}

\begin{solution}
$d\theta = \partial\theta$ is a holomorphic $2$-form, so $d\theta \wedge \overline{d\theta}$ is a form of type $(2,2)$. In local
coordinates, $d\theta = f\,dz_1 \wedge dz_2$, and $d\theta \wedge \overline{d\theta} = |f|^2\,dz_1 \wedge dz_2 \wedge d\bar z_1 \wedge d\bar z_2$,
which is $4|f|^2\,dx_1 \wedge dy_1 \wedge dx_2 \wedge dy_2$, a non-negative multiple of the volume form. On the other
hand $d\theta \wedge \overline{d\theta} = d(\theta \wedge \overline{d\theta})$, because $\overline{d\theta}$ is closed. So its integral
vanishes by Stokes' theorem, which forces $f = 0$. The Iwasawa manifold (Counterexample~\ref{kahler:counterexample-iwasawa})
shows that this fails in dimension $3$.
\end{solution}

\subsection{Hard Lefschetz and the Hodge--Riemann Relations}\label{kahler:subsection-exercises-lefschetz-theorems}

\begin{exercise}\label{kahler:exercise-signatures}
Compute the signature of $\CC\PP^2$, of $\CC\PP^1 \times \CC\PP^1$, and of a smooth quartic surface in $\CC\PP^3$ from
Theorem~\ref{kahler:theorem-hodge-index}, using the Hodge numbers of Section~\ref{kahler:section-hodge-decomposition}.
\end{exercise}

\begin{solution}
For a surface $\sigma = 2 + 2p_g - h^{1,1}$ (the $q$-terms cancel). For $\CC\PP^2$: $p_g = 0$, $h^{1,1} = 1$, so $\sigma = 1$. For
$\CC\PP^1 \times \CC\PP^1$: $p_g = 0$ and $h^{1,1} = 2$ by Exercise~\ref{kahler:exercise-kunneth-hodge}, so $\sigma = 0$, matching the
hyperbolic form of Counterexample~\ref{kahler:counterexample-non-primitive-sign}. For the quartic (Example~\ref{kahler:example-quartic}):
$p_g = 1$ and $h^{1,1} = 20$, so $\sigma = -16$.
\end{solution}

\begin{exercise}\label{kahler:exercise-h1-injective}
Let $X$ be compact K\"ahler of dimension $n \geq 1$ and $0 \neq a \in H^1(X; \RR)$. Show that $a \cup [\omega]^{n-1} \neq 0$, and that
$(a, b) \mapsto \int_Xa \cup b \cup [\omega]^{n-1}$ is a symplectic form on $H^1(X; \RR)$.
\end{exercise}

\begin{solution}
The first statement is hard Lefschetz for $k = 1$. The form is $Q$ for $k = 1$, so it is skew-symmetric and nondegenerate by
Theorem~\ref{kahler:theorem-hodge-riemann}(1). This gives another proof that $b_1$ is even.
\end{solution}

\begin{exercise}\label{kahler:exercise-b2-kahler-surface}
Let $X$ be a compact K\"ahler surface with $b_2(X) = 1$. Show that $p_g = 0$ and $h^{1,1} = 1$, and that a generator $e$ of
$H^2(X; \ZZ)/\mathrm{torsion} \cong \ZZ$ satisfies $\int_Xe^2 = 1$.
\end{exercise}

\begin{solution}
By Theorem~\ref{kahler:theorem-hodge-decomposition}, $b_2 = 2p_g + h^{1,1}$, and $h^{1,1} \geq 1$
(Corollary~\ref{kahler:corollary-hodge-numbers}), so $p_g = 0$ and $h^{1,1} = 1$. By
Corollary~\ref{kahler:corollary-surfaces} the intersection form has signature $(1, 0)$, so $\int e^2 > 0$. By Poincar\'e duality
the cup product pairing on $H^2(X; \ZZ)/\mathrm{torsion}$ is unimodular (Corollary~\ref{poincare-duality:corollary-cup-pairing}), so
$\int e^2 = \pm 1$, hence $1$.
\end{solution}

\begin{exercise}\label{kahler:exercise-primitive-cp1-squared}
Determine the Lefschetz decomposition of $H^\bullet(\CC\PP^1 \times \CC\PP^1; \RR)$ for $[\omega] = h_1 + h_2$, and for
$[\omega] = h_1 + 2h_2$.
\end{exercise}

\begin{solution}
In both cases $P^0 = \RR$ and $H^4 = L^2P^0$. For $[\omega] = h_1 + h_2$, $P^2$ is the kernel of $L \colon H^2 \to H^4$,
$a_1h_1 + a_2h_2 \mapsto (a_1 + a_2)h_1h_2$, which is spanned by $h_1 - h_2$, and $H^2 = LP^0 \oplus P^2$. For $[\omega] = h_1 + 2h_2$, $L$
sends $a_1h_1 + a_2h_2$ to $(2a_1 + a_2)h_1h_2$, so $P^2$ is spanned by $h_1 - 2h_2$. The primitive subspace depends on the K\"ahler class.
\end{solution}

\begin{exercise}\label{kahler:exercise-sl2-weights-cohomology}
Show that the character of the $\mathfrak{sl}_2$-representation $H^\bullet(X; \RR)$ determines the Betti numbers and vice versa, and write
$H^\bullet(\CC\PP^n; \RR)$ and $H^\bullet$ of a compact Riemann surface of genus $g$ as sums of irreducibles $V(m)$.
\end{exercise}

\begin{solution}
The weight $k - n$ has multiplicity $b_k$, so the character is $\sum_kb_kt^{k-n}$. For $\CC\PP^n$, $H^\bullet = V(n)$, spanned by
$1, [\omega], \ldots, [\omega]^n$. For a Riemann surface, $H^\bullet = V(1) \oplus 2g\,V(0)$: the classes $1$ and $[\omega]$ span $V(1)$, and $H^1$
consists of primitive classes of weight $0$.
\end{solution}
